\documentclass[11pt]{amsart}
\usepackage[T1]{fontenc}
\usepackage{lmodern}
\usepackage{amssymb,amsmath}
\usepackage[margin=28mm]{geometry}
\usepackage{xcolor}
\usepackage[colorlinks=true,linkcolor=magenta,citecolor=magenta,urlcolor=magenta]{hyperref}

\newtheorem{theorem}{Theorem}[section]
\newtheorem*{theorema}{Theorem A}
\newtheorem*{theoremb}{Theorem B}
\newtheorem*{theoremc}{Theorem C}
\newtheorem{proposition}[theorem]{Proposition}
\newtheorem{conjecture}[theorem]{Conjecture}
\theoremstyle{definition}
\newtheorem{definition}[theorem]{Definition}

\begin{document}

\title[Standard Model algebra]{Quantum permutations and the Standard Model algebra}

\author{Teo Banica and Claude Opus}

\begin{abstract}
The finite part of the Standard Model, in the Chamseddine--Connes approach, has algebra $A_F=\mathbb C\oplus\mathbb H\oplus M_3(\mathbb C)$, and its quantum isometry groups were computed by Bhowmick, D'Andrea and Dabrowski. We explain how the quantum automorphism group of the complex algebra $B_F=\mathbb C\oplus\mathbb C\oplus M_2(\mathbb C)\oplus M_3(\mathbb C)$ generated by $A_F$, with its canonical trace, is a cocycle twist of the quantum permutation group $S_{15}^+$, and how the color parts of the quantum isometry groups are twists of $S_9^+$, and of a quantum group $G_3\subset S_9^+$ preserving an oriented area tensor on $\mathbb F_3^2$. We prove that $G_3$ is infinite, with classical version of order $216$, while the quantum symmetry group of the unoriented area tensor is classical, of order $432$. With the trace coming from the fermions, the coaction preserves the weak and color summands separately, and the quantum automorphism group becomes a twist of $S_4^+\,\hat{*}\,S_9^+$. We also discuss matrix models and maximality, the fermions and the Heisenberg group of order $27$, the untwisting of color, which leads to color Lie algebras in the sense of Rittenberg, Wyler and Scheunert, the spectral action, and the mixing matrices.
\end{abstract}

\maketitle

\tableofcontents

\section*{Introduction}

This paper is the outline of a program, whose aim is to understand the Standard Model of particle physics in terms of quantum permutations, by using the noncommutative geometry approach of Chamseddine and Connes, and the twisting methods for quantum permutation groups.

The Standard Model of particle physics can be formulated, following Chamseddine and Connes, as the product of an ordinary spin manifold $M$ with a finite noncommutative space $F$, whose algebra is $A_F=\mathbb C\oplus\mathbb H\oplus M_3(\mathbb C)$, see \cite{cc1,cc2,ccm,con,cm}. The gauge group, which is $(U_1\times SU_2\times SU_3)/\mathbb Z_6$, the Higgs field, which appears as an inner fluctuation of the Dirac operator, and the spectral action all emerge from this finite geometry.

It is natural to ask whether $F$ has quantum symmetries, in the sense of Wang \cite{wan}. This question was answered by Bhowmick, D'Andrea and Dabrowski in \cite{bdd}, with the quantum isometry group of $F$ being $C(U_1)*A_{aut}(M_3(\mathbb C))$ in their complex setting, and $C(U_1)*C(PU_3)$ in their real setting, where the coactions must extend to the complexification of the real algebra $A_F$. Here $U_1$ is a circle acting on the weak summand, and the color part is free in the complex setting, and classical in the real setting.

In this paper we revisit this question by using quantum permutations instead of quantum rotations. Our starting point is the twisting result $A_{aut}(\mathbb C_\sigma[G])\simeq A_{aut}(\mathbb C[G])^\sigma$ from \cite{bbc}, which shows that the matrix algebra $M_N(\mathbb C)$ is a twisted version of the algebra of functions on $N^2$ points, and that the quantum group $PU_N^+$ is a twist of $S_{N^2}^+$. This extends to arbitrary finite dimensional $C^*$-algebras, see \cite{behy}: any such algebra $B$, with its canonical trace, is a twisted version of $N=\dim B$ points, in the sense that its quantum automorphism group is a cocycle twist of $S_N^+$. In particular, the complex algebra $B_F=\mathbb C\oplus\mathbb C\oplus M_2(\mathbb C)\oplus M_3(\mathbb C)$ generated by the Chamseddine--Connes algebra corresponds in this way to $15=1+1+4+9$ points. We will present a short proof of all this, which also covers the quantum subgroups defined by extra invariant maps, that we will need.

In this picture the weak sector corresponds to $S_4^+$, which is classical in disguise, untwisting into $SO_3$, and the color sector corresponds to $S_9^+$, which is free on both sides, untwisting into $PU_3^+$. However, all this uses the canonical trace of $B_F$, and with this trace the 15 points get mixed by $S_{15}^+$, and even the electroweak part alone, corresponding to $1+1+4=6$ points, has a free symmetry, coming from $S_6^+$. The physically relevant trace is rather the one coming from the fermions. This trace has unnormalized weights $8,4,4,4$, per generation, on the minimal projections of the 4 summands, and, after normalization, the operator $mm^*$ acts on these summands by the scalars $4,8,16,24$, which are all distinct, and this forces the coaction to preserve each summand. To be more precise, we have the following result:

\begin{theorema}
With the trace coming from its action on the fermions, the quantum automorphism group of $B_F$ is the twist of $S_4^+\,\hat{*}\,S_9^+\subset S_{15}^+$, the quantum subgroup preserving each of the 4 blocks of the partition $15=1+1+4+9$. Equivalently, its algebra of functions is the free product $C(SO_3)*C(PU_3^+)$.
\end{theorema}

Thus, with the physical trace, the coaction preserves the weak and color summands separately, their coefficient algebras appearing in a free product, and freeness appears exactly at the level of quarks. The weak part, which is the twist of $S_4^+$, is then reduced by the Dirac operator to a circle, according to \cite{bdd}, acting on the off-diagonal matrix units of $M_2(\mathbb C)$, while the color part survives intact.

Our main result concerns the real case, and the color part. The classical color group in the real case is characterized by the fact that it preserves the opposite product of $M_3(\mathbb C)$ as well, and on the permutation side the opposite product becomes an oriented area tensor $b$ on the 9 points of $\mathbb F_3^2$. We have the following result:

\begin{theoremb}
The quantum subgroup $G_3\subset S_9^+$ preserving the oriented area tensor on $\mathbb F_3^2$ is a twist of the real color group $PU_3$. It is infinite and noncommutative, and its classical version is the affine symplectic group $\mathbb F_3^2\rtimes SL_2(3)$, of order $216$. By contrast, the quantum symmetry group of the unoriented area tensor $Re(b)$ is classical, equal to $\mathbb F_3^2\rtimes GL_2(3)$, of order $432$.
\end{theoremb}

The first part works in fact for any finite abelian group $H$, with a twist of $PU_N$ appearing inside $S_{N^2}^+$, where $N=|H|$, and the description of this twist by a single oriented triangle tensor is what we would like to emphasize here. As for the second part, the proof uses the fact that the third point on the line through 2 points $a,c$ of $\mathbb F_3^2$ is $-a-c$, and the commutation of the coordinates comes from a general quasigroup argument.

We discuss as well a number of related topics. We construct explicit matrix models for the quantum groups $G_H$, coming from the Weyl operators, we determine exactly when they are commutative, we compute the Hopf image of any family of them, the classical version of $G_H$ corresponding to the Clifford group, and we discuss the maximality of the classical version inside $G_H$, which holds precisely when $H$ has prime exponent. Regarding the fermions, the color triplets cannot be untwisted one at a time, by a determinant argument, and the projective clock and shift cannot be lifted to commuting symmetries of the finite geometry, but the Heisenberg group of order $27$, which is isomorphic to the flavor group $\Delta(27)$, acts on the fermion space of the Standard Model by symmetries of the whole finite geometry. Also, keeping the usual quark module, with the observed connected CKM mixing, at most 3 color points can act. We also discuss some classicality mechanisms, including one coming from the bosonic spectral action: under a coefficientwise reality condition, the natural transformation of 1-forms preserves the bosonic action on every tested geometry precisely for the classical finite quantum groups, while for the actual Standard Model geometry a 2-sided spatial implementer restores the bosonic invariance.

The question of untwisting is in fact the key question raised by the present work. Indeed, the above results show that the algebra of the finite geometry of the Standard Model is, in a precise sense, a twisted version of the algebra of functions on a finite set of points, and the natural problem is that of understanding what happens when the whole theory, including the fermions, the gauge fields and the action, is transported to the untwisted side, where the symmetries are given by quantum permutations. We discuss as well the untwisting of color, at the level of the various objects involved. The twist preserves the dimensions of the representations, and the dimensions of the irreducible representations of the color quantum group $S_9^+$ turn out to be the Fibonacci numbers $1,8,55,377,\ldots$, those of $G_3$ being those of $PU_3$. On the other hand, for $G_3$, which is the twist of $PU_3$, the twist changes the symmetric braiding, the flip becoming a graded flip, with cube roots of unity as coefficients, which is formally reminiscent of the parastatistics of Greenberg \cite{gre}, which was the origin of color. At the level of Lie algebras, we obtain:

\begin{theoremc}
The Lie algebra $gl_3$ is the twist of the color Lie algebra $\mathbb C[\mathbb Z_3^2]$, with bracket $[e_g,e_h]=(1-w^{\det(g,h)})e_{g+h}$, where $w=e^{2\pi i/3}$. Accordingly, in the Weyl basis, the 8 gluon components are labeled by the nonzero elements of $\mathbb F_3^2$, with the pairs of independent labels interacting via trigonometric structure constants.
\end{theoremc}

Here the color Lie algebras are those of Rittenberg, Wyler and Scheunert \cite{rwy,sch}, and the trigonometric structure constants are those of Patera--Zassenhaus and Fairlie--Fletcher--Zachos \cite{ffz,pz}. As for the quarks, these correspond after untwisting to modules over the twisted group algebra of $\mathbb Z_3^2$, in formal analogy with the discrete torsion of Vafa and Witten \cite{vaf,vwi}. In the weak sector the same happens at $N=2$, with the Pauli matrices instead of the clock and shift matrices, and with the cube roots of unity being replaced by signs.

Finally, we discuss the mixing matrices, where the unitary invariant of Connes \cite{co2} reduces in the commuting square case to complex Hadamard matrices, which are well-known to produce quantum permutation groups, and where the subgroups of the weak quantum group $S_4^+$ meet the classical lepton flavor groups $A_4,S_4,A_5$. We also show that the groups of order $216$ and $432$ found above act on the 4 mutually unbiased bases of $\mathbb C^3$ via $A_4$ and $S_4$, with the trimaximal matrix appearing as a transition matrix between these bases, while $G_3$ itself has no nontrivial action on $k\leq8$ points.

The paper is organized as follows: 1 contains the general twisting results, 2 deals with the Chamseddine--Connes algebra, 3 deals with the physical trace and the weak sector, 4 deals with the real structure, 5 discusses matrix models and maximality, 6 discusses the fermions, 7 discusses the untwisting of color, 8 discusses the spectral action, 9 discusses the mixing matrices, 10 discusses the Yang--Mills theory, as a program for later, and 11 contains a list of questions.

\medskip

{\em Acknowledgments.} We would like to thank GPT for a long series of anonymous referee reports, positive and negative, on the present paper. At some point we even considered inviting him as coauthor, but refrained from doing so, due to ethical reasons. It is our hope that he will join us later, in our series of future particle physics investigations.

\section{Twisting finite quantum spaces}

We use Wang's quantum automorphism groups \cite{wan}, in the tracial setting from \cite{ban}. Given a finite dimensional $C^*$-algebra $A$, with a faithful trace $tr:A\to\mathbb C$, we denote by $A_{aut}(A,tr)$ the universal Hopf algebra coacting on $A$ via a unital $*$-algebra morphism $\alpha:A\to A\otimes A_{aut}(A,tr)$, such that $(tr\otimes id)\alpha=tr(.)1$.

More generally, given a family $T$ of linear maps $T_i:A^{\otimes k_i}\to A^{\otimes l_i}$, we denote by $A_{aut}(A,tr,T)$ the quotient of $A_{aut}(A,tr)$ making all the maps $T_i$ colinear. For $A=\mathbb C^N$ with the uniform measure we obtain the algebra $C(S_N^+)$, and for $A=M_N(\mathbb C)$ with the normalized trace we obtain the algebra $C(PU_N^+)$, see \cite{ban,wan}.

Assume that $\Gamma$ is a finite abelian group, and that $A$ is $\Gamma$-graded, with the trace supported in degree $0$. That is, $A=\oplus_{g\in\Gamma}A_g$ with $A_gA_h\subset A_{g+h}$, $A_g^*=A_{-g}$, $1\in A_0$, and $tr(A_g)=0$ for $g\neq0$, the condition $A_g^*=A_{-g}$ being needed for having a grading compatible with the involution. Given a bicharacter $\sigma:\Gamma\times\Gamma\to\mathbb T$, we can then define a twisted product on $A$, as follows:
$$a\cdot b=\sigma(g,h)ab\quad,\quad\forall a\in A_g,b\in A_h$$

We denote by $A_\sigma$ the resulting algebra, and we keep the same trace on it. The grading of $A$ is a coaction of the group algebra $\mathbb C[\Gamma]$, namely $a\to a\otimes g$ for $a\in A_g$, which preserves the trace, so we have a Hopf algebra map $\pi:A_{aut}(A,tr)\to\mathbb C[\Gamma]$. We can therefore twist $A_{aut}(A,tr)$ by using the cocycle $\sigma$, via the map $\pi$, as in \cite{bbc,doi}. We denote the result by $A_{aut}(A,tr)^\sigma$, and we work at the level of the dense Hopf $*$-algebras, the $C^*$-algebras being then obtained by universal completion. The product of the twisted Hopf algebra is given by the following formula, in Sweedler notation:
$$[x][y]=\sum\sigma^{-1}(\pi(x_1),\pi(y_1))\sigma(\pi(x_3),\pi(y_3))[x_2y_2]$$

In a homogeneous basis of $A$, with indices standing for degrees, the coefficients of the coaction $\alpha(e_g)=\sum_ae_a\otimes u_{ag}$ are therefore multiplied according to the formula $[u_{ag}][u_{bh}]=\frac{\sigma(g,h)}{\sigma(a,b)}[u_{ag}u_{bh}]$, the factor $\sigma(a,b)$ coming from the product in the first tensor leg, and the factor $\sigma(g,h)$ coming from the product of the inputs, as in the proof of \cite[Theorem 1.3]{bbc}.

For maps $T:A^{\otimes k}\to A^{\otimes l}$ which preserve the degree, we define $T^\sigma=\Omega_l^{-1}T\Omega_k$, where $\Omega_k$ multiplies $a_1\otimes\ldots\otimes a_k$, with $a_i\in A_{g_i}$, by the scalar $\prod_{j=1}^{k-1}\sigma(g_1+\ldots+g_j,g_{j+1})$. Observe that the multiplication of $A_\sigma$ is $m^\sigma=m\Omega_2$. With these conventions, we have the following result, extending \cite{bbc}:

\begin{theorem}
Assume that $(A,tr)$ is $\Gamma$-graded as above, with the trace supported in degree $0$, and let $T$ be a family of degree preserving maps. We have then
$$A_{aut}(A_\sigma,tr,T^\sigma)\simeq A_{aut}(A,tr,T)^\sigma$$
at the level of the dense Hopf $*$-algebras, and of their universal $C^*$-completions, and in particular these 2 Hopf algebras are isomorphic as coalgebras, via an isomorphism mapping coefficients to coefficients, and preserving the Haar integrals.
\end{theorem}

\begin{proof}
Let us set $L=A_{aut}(A,tr,T)$, and let $\pi:L\to\mathbb C[\Gamma]$ be the Hopf algebra map coming from the grading. As explained in \cite{doi,maj}, the twisting by the 2-cocycle $\sigma(\pi\otimes\pi)$ does not change the coalgebra, and produces a monoidal equivalence between the categories of comodules over $L$ and over $L^\sigma$, which is the identity on the underlying spaces and on the morphisms, and whose tensor structure is given by the maps $\Omega$ above. In other words, a linear map $T:A^{\otimes k}\to A^{\otimes l}$ which preserves the degree is $L$-colinear if and only if $T^\sigma=\Omega_l^{-1}T\Omega_k$ is $L^\sigma$-colinear.

We apply this to the multiplication $m$, to the unit, to the trace, which is supported in degree $0$, so that $tr^\sigma=tr$, and to the maps $T_i$. We conclude that the coaction of $L$ on $A$, regarded as a linear map, is a coaction of $L^\sigma$ on $A_\sigma$, which makes colinear the multiplication $m^\sigma=m\Omega_2$, the unit, the trace and the maps $T_i^\sigma$. Regarding the involutions, the one of $A_\sigma$ is given by $a\to\overline{\sigma(g,-g)}a^*$ for $a\in A_g$, and the compatibility of the coaction with the involutions follows as in \cite{bbc}, the values of $\sigma$ being in $\mathbb T$. Thus, by universality, we obtain a Hopf algebra map $\Phi:A_{aut}(A_\sigma,tr,T^\sigma)\to L^\sigma$, mapping coefficients to coefficients.

The algebra $A_\sigma$ is graded by the same grading, with the trace supported in degree $0$, and the maps $T_i^\sigma$ preserve the degree, so by symmetry, using the cocycle $\sigma^{-1}$, we have as well a Hopf algebra map $\Psi:L\to A_{aut}(A_\sigma,tr,T^\sigma)^{\sigma^{-1}}$, mapping coefficients to coefficients. By universality, $\Phi$ commutes with the maps to $\mathbb C[\Gamma]$ coming from the grading, so it can be twisted by $\sigma^{-1}$, and $\Phi^{\sigma^{-1}}\Psi$ and $\Psi^\sigma\Phi$ are the identity on the coefficients, which generate the algebras. Thus $\Phi$ is an isomorphism. Finally, since the twisting does not change the coalgebra, the Haar integrals coincide on the common vector space, as explained in \cite{bbc}, the products of coefficients appearing in the moments being of course transported via the above formula.
\end{proof}

For related deformation results, in the general setting of $\delta$-forms, see \cite{behy}. In what follows we will only need Theorem 1.1, and its particular cases.

Let us discuss now the basic example. Let $H$ be a finite abelian group, with $|H|=N$, and consider the group $\Gamma=\widehat{H}\times H$. We use the following bicharacters on $\Gamma$:
$$\sigma((\lambda,x),(\mu,y))=\mu(x)\quad,\quad d((\lambda,x),(\mu,y))=\lambda(y)\overline{\mu(x)}$$

Here $d$ is the symplectic form on $\Gamma$, which is nondegenerate and alternating, and we have $\sigma(g,h)d(g,h)=\sigma(h,g)$. We use the group algebra $A=\mathbb C[\Gamma]$, with its standard basis $\{e_g\}_{g\in\Gamma}$, graded by $\deg(e_g)=g$, and with the trace $tr(e_g)=\delta_{g0}$. Via the Fourier transform, $\mathbb C[\Gamma]\simeq C(\widehat{\Gamma})$ is the algebra of functions on $N^2$ points, with the uniform measure. Observe also that the central extension of $\Gamma$ by $\mathbb T$ defined by the cocycle $\sigma$ is the Heisenberg group of $H$, and that $d$ is its commutator form, $d(g,h)=\sigma(h,g)/\sigma(g,h)$. In the main case of interest, $H=\mathbb Z_n$, we have $\Gamma=\mathbb Z_n^2$, and with $w=e^{2\pi i/n}$ our bicharacters are given by $\sigma((i,j),(k,l))=w^{kj}$ and $d((i,j),(k,l))=w^{il-kj}$. On the other hand, the twisted algebra is a matrix algebra, coming from the clock and shift matrices, a construction which goes back to Weyl and Schwinger, and which is well-known in quantum information, where the corresponding operators are called generalized Pauli matrices. To be more precise, we have the following result:

\begin{proposition}
We have an isomorphism $\mathbb C_\sigma[\Gamma]\simeq M_N(\mathbb C)$, mapping the trace to the normalized trace, given by $e_{(\lambda,x)}\to M_\lambda T_x$, where $M_\lambda$ is the multiplication by $\lambda$ on $l^2(H)$, and $T_x$ is the translation $(T_xf)(z)=f(z+x)$. For $H=\mathbb Z_n$, in terms of the matrix units $e_{ab}\in M_n(\mathbb C)$, and with $w=e^{2\pi i/n}$, this isomorphism is given by:
$$e_{(i,j)}\to\sum_kw^{ki}e_{k,k+j}$$
\end{proposition}

\begin{proof}
We have $T_xM_\mu T_x^{-1}=\mu(x)M_\mu$, so $M_\lambda T_xM_\mu T_y=\mu(x)M_{\lambda\mu}T_{x+y}$, which is the twisted product. The operators $M_\lambda T_x$ are linearly independent, and have trace $0$ unless $(\lambda,x)=(1,0)$, so we obtain an isomorphism, mapping the trace to the normalized trace. The formula at $H=\mathbb Z_n$ is a reformulation of this.
\end{proof}

By combining the above results, we recover the following result from \cite{bbc}:
$$A_{aut}(M_N(\mathbb C),tr)\simeq C(S_{N^2}^+)^\sigma$$

At $N=2$ the left term is $C(SO_3)$, see \cite{ban,bddd}, and we recover the fact that $S_4^+$ is a twist of $SO_3$, from \cite{bbi}. At $N\geq3$ both $PU_N^+$ and $S_{N^2}^+$ are infinite and noncommutative. Observe that the twist changes the multiplication of the algebra, and this is unavoidable, due to the following fact:

\begin{proposition}
For $N\geq2$, any trace preserving coaction of $C(S_{N^2}^+)$ on the unchanged algebra $M_N(\mathbb C)$ is trivial.
\end{proposition}

\begin{proof}
The coaction gives an $N^2$-dimensional corepresentation of $S_{N^2}^+$ having a fixed vector, and if nontrivial, this must be the magic representation $1+r_1$, because the other irreducible representations have dimension $>N^2-1$. By restricting to the classical subgroup $S_{N^2}$, a transposition must act on $M_N(\mathbb C)$ by an automorphism of trace $N^2-2$. But an involutive automorphism $Ad_U$ has trace $|Tr(U)|^2=(N-2k)^2$, and $N^2-2$ is never of this form.
\end{proof}

Thus, the multiplication of the algebra must change, and the twisting from Theorem 1.1 is the natural way of changing it. Summarizing, we have obtained the following result:

\begin{theorem}
For any finite abelian group $H$ with $|H|=N\geq2$, the quantum group $PU_N^+$ is the twist of $S_{N^2}^+$ by the cocycle $\sigma$ on $\widehat{H}\times H$, while $S_{N^2}^+$ itself has no nontrivial action on the unchanged algebra $M_N(\mathbb C)$.
\end{theorem}

\begin{proof}
This follows from Theorem 1.1, Proposition 1.2 and Proposition 1.3.
\end{proof}

Summarizing, the matrix algebra $M_N(\mathbb C)$ is a twisted version of the algebra of functions on $N^2$ points, and its quantum symmetries are twisted quantum permutations, with the twisting being unavoidable. In the next section we extend this to arbitrary finite dimensional $C^*$-algebras, and in particular to the Chamseddine--Connes algebra, with its summands $M_2(\mathbb C)$ and $M_3(\mathbb C)$ corresponding to $4$ and $9$ points.

\section{The Chamseddine-Connes algebra}

Before getting into the Standard Model, let us extend the above to arbitrary finite dimensional $C^*$-algebras. Such an algebra is a direct sum of matrix algebras, $B=\oplus_iM_{n_i}(\mathbb C)$, and its canonical trace, coming from the left regular representation, is given by $tr=\oplus_i\frac{n_i^2}{N}tr_i$, where $N=\dim B=\sum_in_i^2$, and where $tr_i$ are the normalized traces, see \cite{ban}. This is the unique trace on $B$ satisfying the $\delta$-form condition $mm^*=N\cdot id$, which is the natural condition for the quantum automorphism groups, as explained in \cite{ban}, and as we will see in section 3 below. We have the following result, which is known, see \cite{behy}:

\begin{theorem}
Any finite dimensional $C^*$-algebra $B=\oplus_iM_{n_i}(\mathbb C)$, with its canonical trace, is a twisted version of $N=\sum_in_i^2$ points, in the sense that
$$A_{aut}(B,tr)\simeq C(S_N^+)^\sigma$$
with $\sigma=\prod_i\sigma_i$ being the product of the cocycles $\sigma_i$ from section 1, on the group $\Gamma=\prod_i\Gamma_i$, with $\Gamma_i=\widehat{H_i}\times H_i$ and $|H_i|=n_i$, acting on the $N$ points by translations on the blocks.
\end{theorem}

\begin{proof}
This is proved in \cite[Theorem 3.1]{behy}, with a construction using translations on the individual blocks and a product cocycle, and here is a short proof, in our setting. We grade the algebra $A=\oplus_i\mathbb C[\Gamma_i]$ by $\Gamma$, with the group algebras graded via the inclusions $\Gamma_i\subset\Gamma$. Via the Fourier transform we have $A\simeq\mathbb C^N$, and the canonical trace of $A$ is the uniform measure, which is supported in degree $0$. By Proposition 1.2, applied on each block, the twisted algebra is $A_\sigma=B$, with the trace becoming the canonical trace of $B$. Thus Theorem 1.1 applies, and gives the result, the map $\pi:C(S_N^+)\to\mathbb C[\Gamma]$ coming from the action of the dual group $\widehat{\Gamma}\simeq\Gamma$ on the $N$ points.
\end{proof}

This is consistent with the results in \cite{ban,dvv}, where $A_{aut}(B,tr)$ was shown to have the fusion rules of $SO_3$, for $N\geq4$, and to be monoidally equivalent to $C(S_N^+)$. Theorem 2.1 is a stronger statement, because a cocycle twist is a monoidal equivalence which preserves the dimensions of the corepresentations.

The finite part of the Standard Model is the real spectral triple $F=(A_F,H_F,D_F,\gamma_F,J_F)$, with algebra $A_F=\mathbb C\oplus\mathbb H\oplus M_3(\mathbb C)$, see \cite{cc2,ccm,con,cm,vsu}. For the general theory of finite spectral triples, see Krajewski \cite{kra} and Paschke--Sitarz \cite{psi}, the latter discussing as well the symmetries coming from Hopf algebras, and for the relation with Morita equivalence, see \cite{dda}. Let us describe $F$ explicitly, following \cite{ccm,vsu}. For one generation, the particles are the following $16$ fermions, with the quarks coming in $3$ colors, and with $R,L$ standing for the right-handed and left-handed components:
$$\begin{array}{|c|c|c|c|c|}
\hline
\text{leptons}&\nu_R&e_R&\nu_L&e_L\\
\hline
\text{quarks}&u_R&d_R&u_L&d_L\\
\hline
\end{array}$$

Thus the space of particles is $H_f=\mathbb C^{16}$, with the quarks contributing $4\times3=12$ dimensions, and the leptons $4$ dimensions. Together with the antiparticles, and with the $3$ generations, we obtain the following Hilbert space, of dimension $96$:
$$H_F=(H_f\oplus\overline{H_f})\otimes\mathbb C^3$$

The algebra $A_F$ acts on $H_F$ as follows, with $a=(\lambda,q,m)$. On the particles, the quaternion $q$, viewed as a $2\times2$ complex matrix, acts on the left-handed doublets $(\nu_L,e_L)$ and $(u_L,d_L)$, while $\lambda$ acts on $\nu_R,u_R$, and $\bar{\lambda}$ acts on $e_R,d_R$. On the antiparticles, $\lambda$ acts on the antileptons, and $m$ acts on the color index of the antiquarks, the leptons being regarded in this way as having a 4th color. The generation space $\mathbb C^3$ is not touched by the algebra.

The real structure is the antiunitary $J_F(\xi,\bar{\eta})=(\eta,\bar{\xi})$, exchanging particles and antiparticles, and the grading $\gamma_F$ is $+1$ on the left-handed particles and the right-handed antiparticles, and $-1$ on the other fermions. These satisfy the KO-dimension $6$ relations, $J_F^2=1$, $J_FD_F=D_FJ_F$, $J_F\gamma_F=-\gamma_FJ_F$. As for the Dirac operator, this is of the following form:
$$D_F=\begin{pmatrix}S&T^*\\ T&\bar{S}\end{pmatrix}$$

Here $S$ acts on the particles, connecting the left and right components, and is built from the Yukawa matrices $Y_\nu,Y_e,Y_u,Y_d\in M_3(\mathbb C)$, acting on the generation space. Up to the normalization coming from the Higgs vacuum, the singular values of these matrices are the Dirac masses of the fermions, and the mixing matrices CKM and PMNS measure the relative position of the diagonalizing bases, for the up and down quarks, and for the neutrinos and electrons, with the light neutrino masses and the PMNS matrix coming in fact from the full neutral mass matrix, including the Majorana term below. As for $T$, this is the Majorana mass term, connecting $\nu_R$ with its antiparticle.

The algebra $A_F$ is real, and the complex algebra that it generates on $H_F$ is the following algebra, of dimension $15$, see \cite{bdd}:
$$B_F=\mathbb C\oplus\mathbb C\oplus M_2(\mathbb C)\oplus M_3(\mathbb C)$$

Here the 2 copies of $\mathbb C$ appear because a complex number $\lambda\in\mathbb C\subset A_F$ acts on some fermions by $\lambda$, and on others by $\bar{\lambda}$, so that the complex span of these operators is $\mathbb C\oplus\mathbb C$, and $M_2(\mathbb C)$ is the complexification of $\mathbb H$. Modulo many things, we have the following dictionary, for the summands of $B_F$:
$$\begin{array}{|c|c|c|c|}
\hline
\text{summand}&\text{gauge group}&\text{interaction}&\text{points}\\
\hline
\mathbb C\oplus\mathbb C&U_1&\text{hypercharge}&1+1\\
\hline
M_2(\mathbb C)&SU_2&\text{weak}&4\\
\hline
M_3(\mathbb C)&SU_3&\text{strong}&9\\
\hline
\end{array}$$

Observe that this dictionary is only a heuristic one, the gauge group of the Standard Model, which acts on the fermions as $(U_1\times SU_2\times SU_3)/\mathbb Z_6$, coming from the unimodular unitary group of $A_F$, and not of $B_F$, with the hypercharges being fixed by the unimodularity condition, see \cite{ccm,vsu}. This gauge group is not to be confused with the quantum symmetry groups discussed below. In what follows, we will be interested in the algebra $B_F$ itself, and in its quantum symmetries. The canonical trace of $B_F$ has weights $\frac{1}{15},\frac{1}{15},\frac{4}{15},\frac{9}{15}$ on the 4 summands, and these are exactly the weights of the uniform measure on $1+1+4+9=15$ points. By applying Theorem 2.1, we obtain the following result:

\begin{theorem}
Consider the group $\Gamma=\mathbb Z_2^2\times\mathbb Z_3^2$, acting on $15=1+1+4+9$ points, with $\mathbb Z_2^2$ acting by translations on the 4 points, and $\mathbb Z_3^2$ acting by translations on the 9 points. We have then
$$A_{aut}(B_F,tr)\simeq C(S_{15}^+)^\sigma$$
with $\sigma=\sigma_2\times\sigma_3$ being the product of the cocycles from section 1, for $H=\mathbb Z_2$ and $H=\mathbb Z_3$.
\end{theorem}

\begin{proof}
This is the particular case of Theorem 2.1, for the algebra $B_F$.
\end{proof}

At the level of the summands, the weak part $M_2(\mathbb C)$ corresponds to $S_4^+$, which untwists into the classical group $SO_3=PU_2$, while the color part $M_3(\mathbb C)$ corresponds to $S_9^+$, which untwists into the free quantum group $PU_3^+$. Thus, summand by summand, freeness appears exactly at the level of quarks. We will see in section 3 that this summand by summand picture is the correct one, when using the trace coming from the fermions.

Note also that the $9$-dimensional magic representation of $S_9^+$ decomposes as $1+8$, and this corresponds, under the twist, to $M_3(\mathbb C)=\mathbb C1\oplus sl_3(\mathbb C)$, with $8$ being the number of gluons. Similarly, $4=1+3$ at the weak level. However, the analogy with the gluons has its limits, due to the following fact:

\begin{proposition}
The $8$-dimensional representations $r_1$ of $S_9^+$ and $sl_3(\mathbb C)$ of $SU_3$ have different cubic invariants, namely:
$$\dim Hom_{S_9^+}(1,r_1^{\otimes3})=1\quad,\quad\dim Hom_{SU_3}(1,sl_3(\mathbb C)^{\otimes3})=2$$
In matrix terms, $PU_3^+$ preserves the trilinear form $Tr(ABC)$, but not $Tr(BAC)$.
\end{proposition}

\begin{proof}
For $S_9^+$ the fixed spaces of the tensor powers of $p=1+r_1$ have dimensions given by the noncrossing partitions, $|NC(k)|=1,1,2,5$ for $k=0,1,2,3$, and by inclusion-exclusion we obtain $5-3\cdot2+3\cdot1-1=1$. For $SU_3$ we have $\dim End_{SU_3}((\mathbb C^3)^{\otimes k})=k!$ for $k\leq3$, and we obtain $6-3\cdot2+3\cdot1-1=2$. The last assertion follows from the fact that $PU_3^+$ preserves the multiplication of $M_3(\mathbb C)$ and the trace, but not the opposite multiplication.
\end{proof}

In familiar notation, the decompositions $8\otimes8=1+8+55$ for $S_9^+$, and $8\otimes8=1+8_s+8_a+10+\overline{10}+27$ for $SU_3$, both of dimension $64$, are distinguished by their invariant tensors. In physics terms, the 2 invariant cubics $Tr(ABC)$ and $Tr(BAC)$ correspond to the symmetric and antisymmetric color tensors $d_{abc}$ and $f_{abc}$, and the free theory only keeps the complex combination $Tr(T_aT_bT_c)=\frac{1}{4}(d_{abc}+if_{abc})$. The structure constants $f_{abc}$ alone come back only with the opposite multiplication, which is related to the real structure, as discussed in section 4 below.

The quantum symmetries of the whole finite geometry $F$, taking into account $H_F$, $D_F$, $\gamma_F$ and $J_F$, were computed by Bhowmick, D'Andrea and Dabrowski in \cite{bdd}, by using Goswami's quantum isometry groups \cite{bgo,gos}, the result being as follows:

\begin{theorem}[Bhowmick, D'Andrea, Dabrowski]
Under the assumptions on the mass matrices from \cite[section 3.1]{bdd}, the quantum group of orientation and real structure preserving isometries of $F$ is
$$QISO^+_J(D_F)=C(U_1)*A_{aut}(M_3(\mathbb C),tr)$$
with $U_1$ acting on $M_2(\mathbb C)$ via $e_{12}\to e_{12}\otimes z$, $e_{21}\to e_{21}\otimes z^*$, and trivially on the rest of $\mathbb C\oplus\mathbb C\oplus M_2(\mathbb C)$. When requiring the coaction to extend to the complexification of the real algebra $A_F$, compatibly with its real form, the result becomes $QISO^+_{\mathbb R}(D_F)=C(U_1)*C(PU_3)$.
\end{theorem}

In the regime from \cite[section 3.1]{bdd}, which requires spectral separation and nonzero simple charged spectra, the effective answer does not depend on the remaining parameters, such as the mixing matrices, while at $D_F=0$ the commutation with $D_F$ imposes no constraint at all. Observe that the weak part $SO_3$ is reduced by $D_F$ to a circle, and that the $U_1$ factor lives on the weak summand. We should mention here that the Higgs field appears in the Chamseddine--Connes approach as an inner fluctuation of the Dirac operator of $M\times F$, while the above concerns the stabilizer of a fixed finite operator $D_F$, so that a relation with the electroweak symmetry breaking would require a field, a vacuum, and the corresponding gauge action. Observe also that the quantum group in Theorem 2.4 is the effective quantum group, acting on the algebra, whose algebra of functions is the Woronowicz subalgebra generated by the coefficients of the adjoint coaction, inside the algebra of a larger, universal quantum group implementing the coaction on $H_F$, and that these 2 objects must not be identified. Also, the global quantum gauge group of \cite{bddd} is different, its mixed weak and color coefficients not being reconstructible from the 2 factors. By combining with Theorem 1.4, we obtain:

\begin{theorem}
In terms of quantum permutations, the result in \cite{bdd} reads
$$QISO^+_J(D_F)\simeq C(U_1)*C(S_9^+)^\sigma$$
in the complex setting, with $\sigma$ being the Heisenberg cocycle on $\mathbb Z_3^2\subset S_9$.
\end{theorem}

\begin{proof}
This follows indeed from Theorem 2.4 and Theorem 1.4, applied at $H=\mathbb Z_3$.
\end{proof}

Observe that the circle $U_1\subset SO_3$ does not contain the Klein subgroup $\mathbb Z_2^2\subset SO_3$ which supports the weak twist from section 1, so the factor $C(U_1)$ has no direct description in terms of $S_4^+$, and remains untouched by the twist.

\section{Physical trace and weak sector}

The results in section 2 use the canonical trace of $B_F$. However, in the spectral setting the algebra comes with its representation on $H_F$, and the natural trace is then the one coming from this representation. In order to discuss this question, let us first record a general fact. Given $B=\oplus_iM_{n_i}(\mathbb C)$ with a faithful trace $tr$, we denote by $t_i$ the trace of a minimal projection of $M_{n_i}(\mathbb C)$, so that $\sum_in_it_i=1$, and we use the scalar product $<a,b>=tr(ab^*)$.

\begin{proposition}
The operator $mm^*$ acts on each summand $M_{n_i}(\mathbb C)$ as the scalar $n_i/t_i$. Thus, any trace preserving coaction on $B$ leaves invariant each of the subalgebras $B_\lambda=\oplus_{n_i/t_i=\lambda}M_{n_i}(\mathbb C)$, with $\alpha(1_\lambda)=1_\lambda\otimes1$.
\end{proposition}

\begin{proof}
The elements $f_{ab}=e_{ab}/\sqrt{t_i}$, with $e_{ab}\in M_{n_i}(\mathbb C)$, form an orthonormal basis of $B$, and we have $m^*(f_{ac})=\frac{1}{\sqrt{t_i}}\sum_bf_{ab}\otimes f_{bc}$, so $mm^*(f_{ac})=\frac{n_i}{t_i}f_{ac}$. A trace preserving coaction is a unitary corepresentation on $(B,<,>)$, as explained in \cite{ban}, so $m^*$ is colinear, and so is $mm^*$, and its eigenspaces $B_\lambda$ are therefore invariant. Finally, $\alpha(1_\lambda)\in B_\lambda\otimes A$ for any $\lambda$, and these elements sum up to $1\otimes1$, so $\alpha(1_\lambda)=1_\lambda\otimes1$.
\end{proof}

For the canonical trace we have $t_i=n_i/N$, and so $mm^*=N$, which is the $\delta$-form condition from \cite{ban}. Regarding now the trace coming from a Hilbert space, the point is that this is invariant under the quantum groups that we are interested in, which are all of Kac type, as follows:

\begin{proposition}
Assume that $B$ acts on a finite dimensional Hilbert space $H$, via a representation $\pi$, and that a coaction of a Kac type quantum group $Q$ on $B$ is implemented by a unitary corepresentation $U$ on $H$, in the sense that $(\pi\otimes id)\alpha(b)=U(\pi(b)\otimes1)U^*$. Then $\alpha$ preserves the trace $Tr_H\circ\pi$.
\end{proposition}

\begin{proof}
With $U=\sum_{ki}e_{ki}\otimes u_{ki}$ and $X=\pi(b)$, we have the following formula:
$$(Tr_H\otimes id)(U(X\otimes1)U^*)=\sum_{ij}X_{ij}\sum_ku_{ki}u_{kj}^*$$

In the Kac case the conjugate corepresentation $\bar{U}=(u_{ki}^*)$ is unitary too, so $\sum_ku_{ki}u_{kj}^*=\delta_{ij}$, and we obtain $Tr_H(X)1$, as claimed.
\end{proof}

Let us compute now the trace $Tr_{H_F}$ on $B_F$. With notations from section 2, the first copy of $\mathbb C$, where $\lambda$ acts, corresponds to $\nu_R,u_R$ and to the 4 antileptons, the second copy, where $\bar{\lambda}$ acts, corresponds to $e_R,d_R$, while $M_2(\mathbb C)$ acts with multiplicity $4$ on the left-handed doublets, and $M_3(\mathbb C)$ acts with multiplicity $4$ on the antiquarks. Thus, for one generation, we have the following table:
$$\begin{array}{|c|c|c|c|}
\hline
\text{summand}&\text{fermions}&Tr_{H_f\oplus\overline{H_f}}(\text{min. proj.})&n_i/t_i\\
\hline
\mathbb C\ (\lambda)&\nu_R,u_R,\bar{\nu}_R,\bar{e}_R,\bar{\nu}_L,\bar{e}_L&8&4\\
\hline
\mathbb C\ (\bar{\lambda})&e_R,d_R&4&8\\
\hline
M_2(\mathbb C)&\nu_L,e_L,u_L,d_L&4&16\\
\hline
M_3(\mathbb C)&\bar{u}_R,\bar{d}_R,\bar{u}_L,\bar{d}_L&4&24\\
\hline
\end{array}$$

Here the quark entries count the 3 colors, the total is $8+4+2\cdot4+3\cdot4=32$, and the last column is computed with the normalized trace. With 3 generations all the traces get multiplied by $3$, and the normalized trace stays the same. We have checked this table with the program {\tt maradona.py}, which builds the representation of $B_F$ on $\mathbb C^{32}$, and computes the operator $mm^*$ for both traces.

The 4 values of $n_i/t_i$ being distinct, Proposition 3.1 shows that the summands of $B_F$ cannot be mixed by quantum symmetries preserving $Tr_{H_F}$. This is in contrast with the canonical trace, where even the electroweak part has genuine quantum symmetries:

\begin{proposition}
The quantum automorphism group of the electroweak algebra $\mathbb C\oplus\mathbb C\oplus M_2(\mathbb C)$, with its canonical trace, is the twist of $S_6^+$, and is infinite and noncommutative.
\end{proposition}

\begin{proof}
The first assertion follows from Theorem 2.1, with $N=6$. Under the twist, the coaction on the algebra corresponds to the magic representation $1+r_1$ of $S_6^+$, with $r_1$ irreducible, so its commutant has dimension $2$. On the other hand, the classical automorphism group $\mathbb Z_2\times SO_3$ leaves invariant the subspaces spanned by $1_1+1_2$, by $1_1-1_2$, by $1_{M_2}$, and $sl_2(\mathbb C)$, so the commutant of its action has dimension at least $4$. Thus our quantum group is not classical, and since its classical version is $\mathbb Z_2\times SO_3$, it is infinite and noncommutative.
\end{proof}

Summarizing, the canonical trace leads to large free quantum groups, mixing all the summands, while the physical trace forces the coaction to preserve each summand. In this latter case, we have the following result, with $S_4^+\,\hat{*}\,S_9^+\subset S_{15}^+$ standing for the quantum subgroup preserving each of the 4 blocks of the partition $15=1+1+4+9$, whose algebra is $C(S_4^+)*C(S_9^+)$:

\begin{theorem}
With the trace $Tr_{H_F}$ coming from the fermions, we have
$$A_{aut}(B_F,Tr_{H_F})\simeq C(S_4^+\,\hat{*}\,S_9^+)^\sigma\simeq C(SO_3)*C(PU_3^+)$$
with $\sigma$ being the cocycle from Theorem 2.2.
\end{theorem}

\begin{proof}
By Proposition 3.1 and the above table, a coaction preserving $Tr_{H_F}$ leaves invariant the 4 summands, and preserves the normalized trace on each of them. Conversely, such a coaction preserves the canonical trace too, and $Tr_{H_F}$. Thus our algebra is $A_{aut}(B_F,tr,T)$, where $tr$ is the canonical trace, and $T$ consists of the maps $\mathbb C\to B_F$ given by $1\to1_i$, which are of degree $0$, so that $T^\sigma=T$. By Theorem 1.1 this is the twist of the quantum permutation algebra of the 15 points preserving the indicator functions of the 4 blocks, which is $C(S_4^+)*C(S_9^+)$. Finally, the universal property of the free product shows that our algebra is also the free product of the algebras $A_{aut}$ of the summands, which are $\mathbb C,\mathbb C,C(SO_3),C(PU_3^+)$.
\end{proof}

This result is in agreement with \cite{bdd}, where the quantum symmetries of the whole finite geometry turn out to be $C(U_1)*A_{aut}(M_3(\mathbb C))$, with a free product between the circle acting on the weak summand and the color part, the weak part $C(SO_3)=C(S_4^+)^\sigma$ being reduced by the Dirac operator to $C(U_1)$. Regarding $S_4^+$ itself, this is the Pauli quantum group from \cite{bcp}, see also section 5. In permutation terms, the picture is that the free product $S_4^+\,\hat{*}\,S_9^+$ gets reduced by $D_F$, with only a circle surviving from the weak part, and with the color quantum group surviving intact.

Let us describe as well the weak part in more detail. At $N=2$ the isomorphism from Proposition 1.2 maps the 4 points of $\mathbb Z_2^2$, or rather the corresponding basis of $\mathbb C[\mathbb Z_2^2]$, to the matrices $1,X,Z,ZX$, where $X,Z$ are the Pauli matrices, and so the magic representation $1+r_1$ of $S_4^+$ corresponds to the decomposition $M_2(\mathbb C)=\mathbb C1\oplus sl_2(\mathbb C)$, with $r_1$ being the adjoint representation of $SO_3$. In physical terms, the 3 nontrivial labels of $\mathbb Z_2^2$, that is, the 3 nontrivial characters of the 4 points, correspond to the 3 weak directions, and the trivial label corresponds to the unit. The twisting is very mild here, because $S_4^+$ and $SO_3$ have the same fusion rules, and the same dimensions of the irreducible representations, namely $1,3,5,7,\ldots$ However, the algebras $C(S_4^+)$ and $C(SO_3)$ are of course very different, the first one being noncommutative, as shown for instance by the Pauli model from section 5, and the passage from one to the other is precisely the twisting. We will come back to the weak sector, and to its untwisting, in section 7 below.

\section{Real structure and area tensors}

In the setting of \cite{bdd}, the real structure $J_F$ is imposed in both cases from Theorem 2.4, and the difference between them comes from the requirement, in the real case, that the coaction extends to the complexification of $A_F$, compatibly with its real form. This complexification has dimension $24$, and contains 2 copies of $M_3(\mathbb C)$, the second one being the conjugate algebra $\overline{M_3(\mathbb C)}$, which is isomorphic to the opposite algebra $M_3(\mathbb C)^{op}$, via $\bar{a}\to a^*$. This suggests that the classical color group should be characterized by the fact that it preserves the opposite product too, and this is indeed the case, as follows:

\begin{proposition}
The quotient of $A_{aut}(M_N(\mathbb C),tr)$ making the opposite product $m^{op}(a\otimes b)=ba$ colinear is $C(PU_N)$.
\end{proposition}

\begin{proof}
Consider the element $F=\sum_{ab}e_{ab}\otimes e_{ba}\in M_N(\mathbb C)\otimes M_N(\mathbb C)$, which is proportional to $m^*(1)$, and is therefore invariant. In the algebra $M_N(\mathbb C)\otimes M_N(\mathbb C)$ the element $F$ is the flip of $\mathbb C^N\otimes\mathbb C^N$, so $F^2=1$ and $F(x\otimes y)F=y\otimes x$. On the other hand, the left and right multiplications by $F$ can be written as follows, with $F$ being inserted in the middle:
$$L_F=(m^{op}\otimes m)(id\otimes F\otimes id)\quad,\quad R_F=(m\otimes m^{op})(id\otimes F\otimes id)$$

Thus, if both $m$ and $m^{op}$ are colinear, so is the flip $\Sigma=R_FL_F$ of $M_N(\mathbb C)\otimes M_N(\mathbb C)$, and this forces the coefficients of the coaction to commute. Thus the quantum group is classical, and so is a closed subgroup of the group of automorphisms of $M_N(\mathbb C)$, which is $PU_N$. Conversely, $PU_N$ preserves $m^{op}$, and this gives the result.
\end{proof}

As an illustration, the program {\tt caniggia.py} checks numerically at $N=2,3$ that the algebra generated by the maps built from $m,m^{op}$ and their adjoints contains $\Sigma$. Let us transport now $m^{op}$ to the permutation side. With notations from section 1, we define a product on $\mathbb C[\Gamma]=C(\widehat{\Gamma})$ by:
$$b(e_g\otimes e_h)=d(g,h)e_{g+h}$$

In the point basis of $\widehat{\Gamma}\simeq\Gamma$, identified via $d$, the map $b$ is given by the following formula, valid for any finite abelian group $H$:
$$b(\delta_a\otimes\delta_c)=\frac{1}{|\Gamma|}\sum_r\overline{d(a-r,c-r)}\,\delta_r$$

Indeed, with $\chi_g(r)=d(g,r)$ we have $\delta_a=|\Gamma|^{-1}\sum_g\overline{d(g,a)}\chi_g$, and inserting this in the formula $b(\chi_g\otimes\chi_h)=d(g,h)\chi_{g+h}$, the summation over $g$ forces $h=a-r$, which gives the result. At $H=\mathbb Z_3$, so that $\Gamma=\mathbb Z_3^2$ and $d(g,h)=w^{\det(g,h)}$ with $w=e^{2\pi i/3}$, we obtain a formula involving the oriented area of the triangle $(r,a,c)$:
$$b(\delta_a\otimes\delta_c)=\frac{1}{9}\sum_rw^{-\det(a-r,c-r)}\delta_r$$

Since $\sigma(g,h)d(g,h)=\sigma(h,g)$, we have $b^\sigma=m_\sigma^{op}$, where $m_\sigma$ is the product of $\mathbb C_\sigma[\Gamma]=M_N(\mathbb C)$. In other words, the twist transports the pointwise product into the product of $M_N(\mathbb C)$, and the area tensor $b$ into the opposite product. This leads to:

\begin{definition}
Given a finite abelian group $H$ with $|H|=N$, we let $G_H\subset S_{N^2}^+$ be the quantum subgroup making the area tensor $b$ colinear, $b\in Hom(p^{\otimes2},p)$, where $p$ is the magic representation. At $H=\mathbb Z_3$ we use the notation $G_3\subset S_9^+$.
\end{definition}

In other words, $G_H$ is the quantum group of quantum permutations of the $N^2$ points which preserve the oriented areas. With this notion in hand, we have the following result:

\begin{theorem}
For any finite abelian group $H$ with $|H|=N\geq2$, the quantum group $G_H\subset S_{N^2}^+$ is a twist of $PU_N$, and is infinite and noncommutative.
\end{theorem}

\begin{proof}
The map $b$ preserves the degree, so Theorem 1.1 applies, and combined with Proposition 4.1 this gives $C(G_H)^\sigma\simeq C(PU_N)$. The algebra $C(G_H)$ is infinite dimensional, because it is isomorphic to $C(PU_N)$ as a coalgebra, while its classical version is finite, being a subgroup of $S_{N^2}$. Thus $G_H$ is infinite and noncommutative.
\end{proof}

Let us mention that embeddings of twists of classical groups into quantum permutation groups, supported on a finite abelian subgroup, go back to \cite[Theorem 8.3]{bbn}, see also \cite{cte}, and that the new point here is the description of the twist of $PU_N$ by a single oriented area tensor.

Regarding now the classical version, this is something quite explicit, because a permutation of the $N^2$ points preserving the oriented areas must be affine. To be more precise, we have the following result:

\begin{proposition}
The classical version of $G_H$ is the affine symplectic group $\Gamma\rtimes Sp(\Gamma,d)$. At $N=2$ we have $G_H=S_4^+$, and at $N=3$ the classical version of $G_3$ is $\mathbb F_3^2\rtimes SL_2(3)$, of order $216$, so that $G_3\neq S_9^+$.
\end{proposition}

\begin{proof}
Identify the $N^2$ points with $\Gamma$, via $d$. A permutation $f$ preserving $b$ preserves the quantities $d(a-r,c-r)$, and up to a translation we can assume $f(0)=0$. Then $d(f(a),f(c))=d(a,c)$ for any $a,c$, and so $d(f(x+y)-f(x)-f(y),f(z))=1$ for any $x,y,z$. Since $f$ is bijective and $d$ is nondegenerate, $f$ is additive, and so is a symplectic automorphism. Conversely, the affine symplectic maps preserve $b$. At $N=2$ the group $PU_2=SO_3$ is the twist of $S_4^+$, so $G_H=S_4^+$. At $N=3$ we have $Sp(\Gamma,d)=SL_2(3)$, and we have also checked by computer, with the program {\tt seedorf.py}, that exactly $216$ of the $9!$ permutations preserve $b$.
\end{proof}

At $N=3$ it is natural to allow orientation reversal too. The program {\tt seedorf.py} finds that exactly $216$ permutations of the 9 points send $b$ to $\bar{b}$, so that exactly $432$ permutations preserve the unoriented area $Re(b)$, forming the group $\mathbb F_3^2\rtimes GL_2(3)$. We will see below that there are no further quantum symmetries of $Re(b)$.

Interestingly, the Heisenberg group of order $27$, which will appear in section 6, is isomorphic to the group $\Delta(27)$, which is well-known as a flavor symmetry group, going back to \cite{bgg}, and whose automorphism group has order $432$ as well, see \cite{hls}, the generalized CP transformations of $\Delta(27)$ models being built from such automorphisms, see \cite{cpv,hls}. However, the natural action of $Aut(\Delta(27))$ on the 9 one-dimensional characters of $\Delta(27)$ fixes the trivial character, so it is not the transitive affine action on the 9 points, and the equality of the orders does not identify these 2 permutation representations. Also, an orientation reversing permutation of the 9 points is not by itself a CP symmetry, which would require compatible transformations of all the multiplets, and of the spacetime and spinor variables, see \cite{cpv}.

Back to the Standard Model, we can now formulate the following result, which is the main result of this paper:

\begin{theorem}
In the real setting, and under the assumptions from \cite[section 3.1]{bdd}, the color part of the quantum isometry group of the finite geometry of the Standard Model, which is $PU_3$ by \cite{bdd}, is the twist of the quantum permutation group $G_3\subset S_9^+$, preserving the oriented area tensor on $\mathbb F_3^2$. Thus, in the real setting, we have
$$QISO^+_{\mathbb R}(D_F)\simeq C(U_1)*C(G_3)^\sigma$$
with $G_3$ being infinite and noncommutative, and with classical version of order $216$.
\end{theorem}

\begin{proof}
This follows from Theorem 2.4, Theorem 4.3 and Proposition 4.4, applied at $H=\mathbb Z_3$.
\end{proof}

Thus, in the permutation picture, the color symmetry of the real Standard Model, which is classical in \cite{bdd}, corresponds to a genuinely quantum permutation group. The same phenomenon appears at $N=2$, where the classical $SO_3$ untwists into $S_4^+$. For a related study of the twists of $SU_3$, see \cite{by}.

Geometrically, the 9 points of $\mathbb F_3^2$ on which $G_3$ acts can be thought of as the 9 inflection points of a plane cubic, the 12 affine lines being the lines joining them, and the classical version of $G_3$ being the Hessian group of order 216, see \cite{ado}. The oriented area tensor $b$ is then an invariant of this configuration, and $G_3$ can be regarded as the quantum symmetry group of the Hesse configuration, equipped with its oriented areas. We will come back to this point of view in section 9, in connection with the mixing matrices.

Let us discuss now the unoriented version of the area tensor. We consider the entrywise real part $Re(b)$ of the tensor $b$, written in the point basis as above, which is invariant under the orientation reversing affine maps too. Observe that taking adjoints in an intertwining relation reverses the order of the products of magic entries, while the entrywise conjugation in the point basis does not follow from such an operation, so preserving $b$ does not imply preserving $Re(b)$. We will need the following general fact:

\begin{proposition}
For a finite quasigroup $(X,*)$, the quantum subgroup of $S_X^+$ making colinear the map $L(v_a\otimes v_c)=v_{a*c}$ is classical, namely the group of automorphisms of $(X,*)$.
\end{proposition}

\begin{proof}
For a magic unitary $p$, the colinearity of $L$ reads $\sum_{x*y=k}p_{xi}p_{yj}=p_{k,i*j}$. Fix $a,b,i,j$, and set $k=a*b$, $P=p_{ai}$, $Q=p_{bj}$, $R=p_{k,i*j}$. By multiplying this relation on the left by $P$, and separately on the right by $Q$, and using the orthogonality on the columns of $p$, and the unique division in the quasigroup, we obtain $PQ=PR=RQ$. Thus $(PQ)^2=PRRQ=PRQ=PQ$, so $PQ$ is a contractive idempotent, and so an orthogonal projection, and so $PQ=(PQ)^*=QP$. Thus the coordinates commute, and the classical points are the permutations preserving the operation $*$.
\end{proof}

For a related rigidity result, in the associative case, see \cite{ksw}. Back now to $\mathbb F_3^2$, the point is that the passage from $b$ to $Re(b)$ kills all the quantum symmetries, as follows:

\begin{theorem}
The quantum subgroup of $S_9^+$ preserving the unoriented area tensor $Re(b)$ on $\mathbb F_3^2$ is the classical group $\mathbb F_3^2\rtimes GL_2(3)$, of order $432$. In particular, this quantum group does not contain $G_3$, and its intersection with $G_3$ is the classical version of $G_3$.
\end{theorem}

\begin{proof}
Set $v_a=\delta_a$ and $1=\sum_rv_r$, and let $C(v_a\otimes v_c)$ be the sum of the vectors $v_r$ with $\det(a-r,c-r)=0$. Since the real parts of the nontrivial cube roots of unity are $-1/2$, we have $Re(b)(v_a\otimes v_c)=-\frac{1}{18}1+\frac{1}{6}C(v_a\otimes v_c)$. For $a\neq c$ the points $r$ on the line through $a,c$ are $a,c,-a-c$, and we obtain, for any $a,c$:
$$C(v_a\otimes v_c)=v_a+v_c+v_{-a-c}+\delta_{ac}(1-3v_a)$$

All the terms here, except for $L(v_a\otimes v_c)=v_{-a-c}$, are built from the unit, the summation functional, the pointwise multiplication and the identity, so they are intertwiners for any magic unitary, and preserving $Re(b)$ is equivalent to preserving $L$. Now $a*c=-a-c$ is a quasigroup operation, so by Proposition 4.6 our quantum group is classical. Now a permutation $f$ preserving $L$ and fixing $0$ satisfies $f(-a)=-f(a)$, and then $f(a+b)=f(a)+f(b)$, so the permutations preserving $L$ are the affine maps. Finally, the intersection with $G_3$ is classical, and so contained in the classical version of $G_3$, which is contained in $\mathbb F_3^2\rtimes GL_2(3)$. The program {\tt gullit.py} checks the above formulae, finds that $432$ permutations preserve $L$, and checks that the Weyl models from section 5 below intertwine $b$, but not $Re(b)$.
\end{proof}

In fact, the orientation of the areas, which is what makes room for quantum symmetries here, is invisible on the incidence structure of the 12 affine lines, and this structure has only classical symmetries. To be more precise, with the quantum automorphism group of a graph being defined by a magic unitary commuting with its adjacency matrix, we have:

\begin{theorem}
The quantum automorphism group of the incidence graph of the Hesse configuration, having as vertices the 9 points of $\mathbb F_3^2$ and the 12 affine lines, is the classical group $\mathbb F_3^2\rtimes GL_2(3)$, of order $432$.
\end{theorem}

\begin{proof}
The degree vector, which is the image of the all-one vector by the adjacency matrix, is fixed, and since its entries are $4$ on the points and $3$ on the lines, the magic unitary is block diagonal, with blocks $u$ on the points and $v$ on the lines. With $I:\mathbb C^{12}\to\mathbb C^9$ being the incidence matrix, we have $uI=Iv$, and $I^*$ intertwines too. Now let $T=Im_l(I^*\otimes I^*)$, with $m_l$ being the pointwise product of the functions on the lines. Since there is a unique line through 2 distinct points, and the 4 lines through a point $a$ cover the other points once, and $a$ itself four times, we have:
$$T(v_a\otimes v_c)=v_a+v_c+v_{-a-c}+\delta_{ac}1$$

Thus, after subtracting universal intertwiners, we obtain the map $L$ from the proof of Theorem 4.7, and by Proposition 4.6 the coefficients of $u$ commute. Also, for $a\neq c$, the vector $m_l(I^*v_a\otimes I^*v_c)$ is the indicator of the line through $a,c$, and these vectors span the functions on the lines, so the coefficients of $v$ lie in the algebra generated by those of $u$, and commute too. Finally, the classical automorphisms are the permutations of the points preserving the lines, which are the affine maps. The program {\tt haaland.py} checks the above formula on all 81 pairs of points, and finds $432$ permutations preserving the lines.
\end{proof}

\section{Matrix models and maximality}

The quantum groups $G_H$ have explicit matrix models, coming from the Weyl operators. Let us denote by $W_g=M_\lambda T_x$, with $g=(\lambda,x)\in\Gamma$, the Weyl operators from Proposition 1.2, which form an orthonormal basis of $M_N(\mathbb C)$, with respect to the scalar product $<a,b>=tr(ab^*)$. Given a unitary $x\in U_N$, we can consider the following matrix of rank one projections, acting on the Hilbert space $M_N(\mathbb C)$:
$$P^x_{rs}=Proj(W_rxW_s^*)$$

Here the indices $r,s$ range over $\Gamma$, identified with the $N^2$ points via $d$. We have the following result, in the spirit of the flat models from \cite{bne}:

\begin{proposition}
The matrix $P^x$ is magic, and the area tensor $b$ intertwines it, so we have a representation $\pi_x:C(G_H)\to B(M_N(\mathbb C))$, given by $p_{rs}\to P^x_{rs}$.
\end{proposition}

\begin{proof}
We have the following formulae, with the scalar product $<a,b>=tr(ab^*)$:
$$<W_rxW_s^*,W_rxW_t^*>=tr(W_s^*W_t)=\delta_{st}\quad,\quad<W_rxW_s^*,W_qxW_s^*>=tr(W_rW_q^*)=\delta_{rq}$$

Thus the rows and columns of $P^x$ consist of orthonormal bases, and $P^x$ is magic. Regarding $b$, consider the operators $V_{ab}(Y)=W_aYW_b^*$, which form a basis of $B(M_N(\mathbb C))$, and the coefficients $t_{ab}(x)=tr(W_a^*xW_bx^*)$ of the conjugation action of $PU_N$ on $M_N(\mathbb C)$, in the Weyl basis. By expanding the rank one projections, and using $W_rW_aW_r^*=d(a,r)W_a$, we obtain, with $m=|\Gamma|$:
$$P^x_{rs}=\frac{1}{m}\sum_{a,b}d(a,r)\overline{d(b,s)}\,t_{ab}(x)V_{ab}$$

Since $V_{ab}V_{cd}=\sigma(a,c)\overline{\sigma(b,d)}V_{a+c,b+d}$, this shows that $p_{rs}\to P^x_{rs}$ is the evaluation at $x$ of the twisted version of the algebra of functions on $PU_N$, which is $C(G_H)$, via the map $f\to\sum_{ab}f_{ab}(x)V_{ab}$, where $f_{ab}$ are the components of $f$ under the left and right translations by the Weyl group, as in \cite[Theorem 4.11]{cte}. Thus $\pi_x$ is well-defined, and $b$ intertwines $P^x$. We have verified the above formula with the program {\tt koeman.py}, for $H=\mathbb Z_3,\mathbb Z_4,\mathbb Z_2^2$, and the intertwining property at $N=3$ with the program {\tt kempes.py}.
\end{proof}

Let us discuss now the commutativity of these models. The Weyl group $K\simeq\Gamma$ sits inside $PU_N$, via $g\to[W_g]$, and its normalizer $C_H$, which is the projective Clifford group, fits into an exact sequence
$$1\to K\to C_H\to Sp(\Gamma,d)\to1$$
with the quotient map being given by the conjugation action on the labels of the Weyl operators, see \cite{gal}. In particular $C_H$ is finite, of order $|\Gamma|\cdot|Sp(\Gamma,d)|$, and we have:

\begin{proposition}
The image of $\pi_x$ is commutative if and only if $[x]\in C_H$. In particular, $\pi_x$ has noncommutative image for any $x$ outside a finite set of points of $PU_N$.
\end{proposition}

\begin{proof}
The column $s=0$ of $P^x$ consists of the projections onto the orthonormal basis $\{W_rx\}_r$. If all the entries of $P^x$ commute, each rank one projection $P^x_{rs}$ commutes with these, so it is diagonal in this basis, and so $xW_s^*$ must be proportional to some $W_rx$, for any $s$. Thus $[x]$ normalizes $K$. Conversely, if $[x]$ normalizes $K$, then each $W_rxW_s^*$ is proportional to some $W_tx$, so all the projections $P^x_{rs}$ are diagonal in the basis $\{W_tx\}_t$, and so commute. The program {\tt benzema.py} checks at $N=3$ that $C_3$ has order $216$, that $\pi_x$ is commutative for $x\in C_3$, and that it is noncommutative for random $x$, and for small perturbations of elements of $C_3$.
\end{proof}

At $N=2$ the Weyl operators are, up to phases, the Pauli matrices, and the models $\pi_x$ with $x\in SU_2$ are the Pauli representation of $C(S_4^+)$, which was shown to be faithful in \cite{bcp}. Thus Proposition 5.1 can be regarded as a generalization of the Pauli representation. In general, the integrated model $x\to\pi_x$, with $x\in PU_N$, is faithful on the dense Hopf $*$-algebra, as shown in \cite[Theorem 4.11]{cte}, where the Hopf image of the full Weyl model is identified with the twist of $PU_N$. More generally, we have the following result, where the Hopf image of a family of representations is the smallest Hopf algebra quotient through which they all factor:

\begin{theorem}
Given a nonempty subset $X\subset PU_N$, the joint Hopf image of the models $\pi_x$ with $x\in X$ is the twist of the algebra of functions on the closed subgroup $\overline{\langle K,X\rangle}\subset PU_N$.
\end{theorem}

\begin{proof}
We work at the level of the dense Hopf $*$-algebras, and we identify the polynomial functions on the twist of $PU_N$ with those on $PU_N$, as vector spaces. Decompose such a function under the left and right translations by $K$, as $f=\sum_{a,b}f_{ab}$. By the proof of Proposition 5.1, see also \cite[Theorem 4.11]{cte}, $\pi_x$ maps $f_{ab}$ to $f_{ab}(x)V_{ab}$, where $V_{ab}(Y)=W_aYW_b^*$, and since these operators form a basis of $B(M_N(\mathbb C))$, finite Fourier inversion shows that $\pi_x(f)=0$ precisely when $f$ vanishes on $KxK$. Similarly, the iterated coproduct being given by evaluation at products, the tensor $(\pi_{x_1}\otimes\ldots\otimes\pi_{x_m})\Delta^{(m-1)}(f)$ vanishes precisely when $f$ vanishes on $Kx_1K\ldots Kx_mK$. The union of these sets, over all words in $X$, is a semigroup, whose closure is a closed subsemigroup of a compact group, and so a closed subgroup, which contains $X$, and $K=(Kx)x^{-1}$ for $x\in X$, so it is $H=\overline{\langle K,X\rangle}$. Thus the largest Hopf ideal contained in the joint kernel of our models is the ideal of functions vanishing on $H$. Conversely, since $K\subset H$, the restriction to $H$ is a Hopf algebra quotient of the twisted algebra, and every model factors through it. We have checked at $N=3$, with the program {\tt benzema.py}, that the kernels of $f\to\pi_x(f)$ and of $f\to f_{|KxK}$ coincide on the coefficients of the magic representation.
\end{proof}

In particular, with $X=C_H$, we obtain that the twist of $C(C_H)$ is a quotient of $C(G_H)$, which is commutative by Proposition 5.2, and has dimension $|C_H|=|\Gamma|\cdot|Sp(\Gamma,d)|$, so it is the algebra of functions on the classical version of $G_H$, computed in Proposition 4.4. Thus the classical version of $G_H$ corresponds, under the twist, to the projective Clifford group $C_H$, which answers a question raised in \cite[section 4]{cte}.

Also, for Haar almost every $x\in PU_N$, the single model $\pi_x$ has Hopf image $G_H$. Indeed, consider the closed subgroup $L=\overline{\langle K,x\rangle}$. The complexified Lie algebra of $L$ is invariant under $Ad(K)$, so it is a sum of Weyl lines $\mathbb CW_g$, and there are finitely many such sums. A proper nonzero such sum has a proper closed stabilizer in $PU_N$, because the adjoint representation is irreducible, and so $x$ must belong to a Haar null set. Also, if $L$ is finite, then $x$ is a torsion element, and these elements form a null set. Thus, outside a null set we have $L=PU_N$, and Theorem 5.3 gives the result.

The restriction to abelian groups $H$ in all this matters. Indeed, consider a nice error basis $\{U_g\}$ indexed by a finite nonabelian group $F$, and restrict the above model to the parameters $x=U_k$. Each vector $U_iU_kU_j^*$ is proportional to another basis element, with label $ikj^{-1}$, so the projections are diagonal in the basis $\{U_g\}$, and so commute. Thus the model factors through evaluations at the permutations $i\to g^{-1}ik$, and the Hopf image is the classical group of left and right translations $(F\times F)/\{(z,z)|z\in Z(F)\}$, of order $|F|^2/|Z(F)|$, while a twist of $C(F)$ would have dimension $|F|$. For instance the 16 operators $Z^aX^bJ^s\otimes Z^sX^t$ with $a,b,s,t\in\mathbb F_2$, where $Z,X$ are the Pauli matrices and $J$ is the Hadamard matrix of order 2, form such a basis of $M_4(\mathbb C)$, with $|F|=16$ and $|Z(F)|=4$, so the Hopf image has order $64$. We have checked all this with the program {\tt thuram.py}, which verifies that these operators form a nice error basis, computes the group law of the labels and its center, and computes the group generated by the above permutations.

Let us discuss now maximality. Under the twist, by \cite[Lemma 4.13]{cte}, the closed quantum subgroups of $G_H$ containing $K$ correspond to the closed subgroups of $PU_N$ containing $K$, and the classical version $\Gamma\rtimes Sp(\Gamma,d)$ corresponds to $C_H$, as explained above. We have the following result, in the spirit of the maximality questions for $S_N\subset S_N^+$, see \cite{cte}:

\begin{theorem}
For a nontrivial finite abelian group $H$, the classical version $\Gamma\rtimes Sp(\Gamma,d)$ is maximal among the closed quantum subgroups of $G_H$ precisely when $H$ has prime exponent, that is, $H=\mathbb Z_p^r$ with $p$ prime.
\end{theorem}

\begin{proof}
By untwisting, the question is whether $C_H$ is maximal in $PU_N$. For $H=\mathbb Z_p^r$, the conjugation by $K$ has distinct $1$-dimensional weight spaces $\mathbb CW_g$ with $g\neq0$ on $sl_N(\mathbb C)$, and the symplectic group acts transitively on the nonzero labels, so $C_H$ acts irreducibly on $sl_N(\mathbb C)$. Thus there is no proper closed overgroup of positive dimension, because the complexified Lie algebra of its connected component would be a nonzero invariant subspace, hence equal to $sl_N(\mathbb C)$, and the overgroup would be $PU_N$ itself. The finite overgroups are excluded by the results of Nebe, Rains and Sloane on the maximality of the complex Clifford groups associated to $\mathbb Z_p^r$, see \cite[Theorems 6.5 and 7.3]{nrs} and \cite{nr2}, via a finite lift: the preimage in $SU_N$ of a finite overgroup of $C_H$ is finite, and adjoining the finite group of scalar phases of their Clifford group, we obtain a finite overgroup of it. Conversely, if the exponent of $H$ is composite, pick a prime $p$ dividing it, and set $D=\Gamma[p]$, which is a proper nonzero characteristic subgroup. The algebra spanned by the operators $W_g$ with $g\in D$ is a proper nonscalar subalgebra of $M_N(\mathbb C)$, and its normalizer in $PU_N$ contains $C_H$, because $C_H$ acts on the labels by symplectic automorphisms, which preserve $D$. This normalizer is proper, because a subalgebra invariant under all the unitary conjugations must be trivial or the whole algebra, and has positive dimension, because it contains the conjugations by the unitaries of our nonscalar subalgebra. By twisting, we obtain an intermediate quantum group.
\end{proof}

For instance at $H=\mathbb Z_9$ the subgroup $D=\Gamma[3]\simeq\mathbb Z_3^2$ is isotropic for $d$, so the algebra spanned by the operators $W_g$ with $g\in D$ is commutative, of dimension $9$, and so is a maximal abelian subalgebra of $M_9(\mathbb C)$, generated by the cubes $X^3,Z^3$ of the shift and clock matrices. Its normalizer in $PU_9$ is the corresponding monomial group, namely a torus extended by $S_9$, which gives an intermediate group. On the other hand, for $H=\mathbb Z_3^2$, which has the same number $81$ of points, there are no intermediate quantum groups. Thus the maximality question depends on the group structure of $H$, and not only on its order.

\section{Fermions and the Heisenberg group}

The fermions of the Standard Model live in $H_F$, which is a module over $A_F$, and the quarks come in color triplets $\mathbb C^3$, which are irreducible modules over $M_3(\mathbb C)$. A natural question is whether the twisting from the previous sections can be performed at the level of the fermions too. The answer is negative, for single triplets:

\begin{proposition}
The following hold:
\begin{enumerate}
\item If $V=\mathbb C^3$ is a color triplet, with $End(V)=M_3(\mathbb C)$ as algebra object, there is no tensor functor extending the inverse twist which sends $End(V)$ to $\mathbb C^9$, with its pointwise product.

\item For $H=\mathbb Z_N$, if the projective Weyl translations of $\mathbb C^N$ lift to commuting unitaries of the form $M_\lambda\otimes A$, $T_x\otimes B$ on $\mathbb C^N\otimes E$, then $N|\dim E$.
\end{enumerate}
\end{proposition}

\begin{proof}
(1) A tensor functor preserves duals and contractions, so it maps the algebra $End(V)=V\otimes V^*$ to $End(F(V))$, whose center is $1$-dimensional, while $\mathbb C^9$ has a $9$-dimensional center.

(2) From $T_xM_\lambda=wM_\lambda T_x$ with $w$ a primitive root of unity, the commutation of $M_\lambda\otimes A$ and $T_x\otimes B$ gives $AB=wBA$, and taking determinants gives $w^{\dim E}=1$.
\end{proof}

Thus the quarks, taken one triplet at a time, do not survive the passage to 9 points. However, the obstruction in (2) disappears in 2 natural ways. First, the Heisenberg group of order 27, which is the central extension of $\mathbb Z_3^2$ generated by clock and shift, acts genuinely on $\mathbb C^3$. Second, the color space can be paired with a second 3-dimensional space, as in $\mathbb C^3\otimes\overline{\mathbb C}^3$, on which the Weyl translations commute, with all 9 characters as weights. In the Standard Model, a natural candidate for this second space would be the space of the 3 generations, but this is excluded, at least under natural requirements, as follows:

\begin{proposition}
Assume that $D_F^2$, restricted to the antiquarks of a given chirality, has an eigenspace which is a single color triplet, as is the case for nondegenerate quark masses. Then there are no commuting unitaries $U,V$ on $H_F$, commuting with $D_F,\gamma_F$, and implementing on $M_3(\mathbb C)\subset A_F$ the adjoint actions of the clock and shift matrices.
\end{proposition}

\begin{proof}
Such unitaries preserve the support of $M_3(\mathbb C)$, which is the space of antiquarks, and commute with $\gamma_F$ and $D_F^2$, so they preserve the above eigenspace $E\simeq\mathbb C^3$. Since $M_3(\mathbb C)$ acts irreducibly on $E$, by Schur's lemma the restrictions of $U,V$ to $E$ are scalar multiples of the clock and shift matrices. But these do not commute, their commutator being a nontrivial cube root of unity.
\end{proof}

Thus the commuting lifts are excluded, even when using the 3 generations, and the genuinely nonabelian Heisenberg group is the right object.

On the other hand, there are real spectral triples with genuine $S_N^+$ symmetry, provided that the fermions are changed. With $L=l^2(X)$, $|X|=N\geq4$, and $P$ the projection onto the constant functions, one can take $H=L\otimes L\otimes\mathbb C^2$, with $C(X)$ acting on the first leg, $D=(P\otimes1+1\otimes P)\otimes\sigma_x$, $\gamma=1\otimes1\otimes\sigma_z$, and $J$ the flip of the 2 copies of $L$, composed with $\sigma_x$ and complex conjugation. This is a real, even spectral triple of KO-dimension $6$, satisfying the first order condition, the commutators $[D,a]$ acting on the first leg, and the opposite algebra acting on the second leg. The quantum group $S_N^+$ acts on it via $U=u_{13}u_{23}$, which commutes with $D$ because the magic unitary $u$ fixes the constant functions, and which implements the coaction on the first leg. As for the real structure, the flip of the 2 legs, combined with the adjoint on the coefficients, leaves $U$ invariant, because $(u_{ij}u_{kl})^*=u_{kl}u_{ij}$, which is a reality condition of coefficientwise type, in the sense of section 8. We have checked all this numerically at $N=4$, with the Pauli model of $S_4^+$ from section 5, with the program {\tt baresi.py}. Thus the obstruction to quantum symmetry comes from the fermion content of the Standard Model, and not from the axioms themselves.

For the Standard Model itself, the color structure of $H_F$ makes the Heisenberg group act in a very natural way, and we have the following result, which shows that this group is the right object:

\begin{theorem}
The Heisenberg group of order $27$, which is isomorphic to the flavor group $\Delta(27)$, acts on $H_F$ by unitaries commuting with $D_F,\gamma_F,J_F$, implementing on $M_3(\mathbb C)\subset A_F$ the adjoint action of the clock and shift matrices, and acting trivially on $\mathbb C\oplus\mathbb H$.
\end{theorem}

\begin{proof}
In the Standard Model, the color index of each quark or antiquark is carried by a tensor factor $\mathbb C^3$, the leptons being colorless, and $D_F$ and $\gamma_F$ act trivially on this factor, because the Yukawa couplings and the Majorana mass term are color-blind. We let the Heisenberg group act via clock and shift on the color factor of the antiquarks, on which $M_3(\mathbb C)$ acts, via the complex conjugate matrices on the color factor of the quarks, and trivially on the leptons. These unitaries commute with $D_F,\gamma_F$, and with the antiunitary $J_F$, which exchanges particles and antiparticles, and they implement the adjoint action of clock and shift on $M_3(\mathbb C)$.
\end{proof}

This action exists for any choice of the Yukawa couplings and of the mixing matrices, and puts no constraint on them. Observe also that, in the flavor models, the group $\Delta(27)$ acts on the generations, and not on color.

Theorem 6.3 suggests performing the twisting at the level of $F$ by a deformation of isospectral type, in the spirit of Connes and Landi \cite{cla}, with the Heisenberg group instead of a torus. However, this cannot work in the usual monoidal framework: the color triplet is an object of the representation category of the Heisenberg group, its endomorphism algebra is an algebra object of this category, and any tensor functor maps it to a full matrix algebra, as in the proof of Proposition 6.1 (1), so the color algebra cannot become $\mathbb C^9$ in this way. Thus, the passage to 9 points at the level of the fermions requires changing the module data. In fact, keeping the usual quark module, such a change cannot produce many color points, as follows:

\begin{proposition}
Keep $H_F$, the weak action, $J_F$, $\gamma_F$ and $D_F$, with invertible quark mass matrices having simple singular spectra, and replace the action of $M_3(\mathbb C)$ by a faithful representation $\rho$ of $\mathbb C^m$ on the antiquarks, with the same support, and commuting with $\gamma_F$. If the order zero and first order conditions hold, then $m\leq3c$, where $c$ is the number of connected components of the bipartite support graph of the CKM matrix, having the up and down mass eigenstates as its 2 vertex sets. In particular, for connected CKM mixing, $m\leq3$.
\end{proposition}

\begin{proof}
For a minimal projection $f\in\mathbb C^m$, set $Q=J_F\rho(f)J_F^{-1}$, acting on the quarks. By order zero, $Q$ commutes with the weak algebra, so on the left quarks we have $Q_L=1_2\otimes T$, with $T$ acting on color and generations. With $M=diag(1_3\otimes Y_u,1_3\otimes Y_d)$ being the mass arrow from the right quarks to the left quarks, the first order condition $[[D_F,e_w],Q]=0$, with $e_w$ being the unit of $\mathbb H$, gives $Q_LM=MQ_R$, and so $[Q_L,MM^*]=0$. Thus $T$ commutes with $1_3\otimes Y_uY_u^*$ and $1_3\otimes Y_dY_d^*$, and the simple spectra show that $T\in M_3(\mathbb C)\otimes\mathbb C^c$, because a matrix which is diagonal in both eigenbases has diagonal entries satisfying $\lambda_i=\mu_j$ when $V_{ij}\neq0$, where $V$ is the CKM matrix. Since $M$ is invertible, each $T$ is nonzero, and since these $T$ are mutually orthogonal projections, there are at most $3c$ of them. The program {\tt aguero.py} checks the commutant computation, for connected, block and monomial CKM matrices.
\end{proof}

Since the observed CKM matrix is connected, the usual quark module allows at most 3 color points, the bound being attained by the diagonal color algebra, and so 9 color points require a substantial change of the module data.

Let us record as well the decomposition of $H_F$ under this action. The center $\mathbb Z_3$ of the Heisenberg group acts by mutually conjugate primitive cube roots of unity on the quarks and on the antiquarks, and trivially on the leptons, and this is the action of the center of $SU_3$, also known as triality. Since the 3-dimensional irreducible representations of the Heisenberg group are determined by their central character, we obtain, for the 3 generations:
$$H_F=24\cdot1\oplus12\cdot\mathbf 3\oplus12\cdot\bar{\mathbf 3}$$

Here the trivial part corresponds to the leptons, and $\mathbf 3,\bar{\mathbf 3}$ are the 3-dimensional representations given by clock and shift, and by their conjugates. Observe that the color singlets, such as those appearing in $qqq$ and $q\bar{q}$, have trivial central character, the converse being false, since for instance $\mathbf 3\otimes\bar{\mathbf 3}=\mathbf 1\oplus\mathbf 8$ for $SU_3$, with both components having trivial central character. We have checked this decomposition, for one generation, with the program {\tt maradona.py}.

\section{Untwisting and color Lie algebras}

We discuss here the untwisting of the color sector, from $M_3(\mathbb C)$ to the $9$ points of $\mathbb F_3^2$. According to Theorem 1.1, the twist preserves the coalgebra structure, and in particular the dimensions of the representations, the fusion rules and the Haar integrals. Thus, all the quantities which can be computed from the representation theory alone are the same on both sides.

As a first illustration, let us compute the dimensions of the irreducible representations of $S_N^+$. For $N\geq4$ these are indexed by $k\in\mathbb N$, with fusion rules $r_1\otimes r_k=r_{k-1}+r_k+r_{k+1}$, and $\dim r_1=N-1$, see \cite{ban}. Thus the dimensions are given by:
$$d_0=1\quad,\quad d_1=N-1\quad,\quad d_{k+1}=(N-2)d_k-d_{k-1}$$

These dimensions should not be confused with the Catalan numbers $C_k=|NC(k)|$, which count the fixed points of the tensor powers of the magic representation, as in Proposition 2.3, and which do not depend on $N$. Here are the first values of these dimensions, computed with the program {\tt rijkaard.py}:
$$\begin{array}{|c|c|c|c|c|c|}
\hline
N&d_1&d_2&d_3&d_4&d_5\\
\hline
4&3&5&7&9&11\\
\hline
5&4&11&29&76&199\\
\hline
6&5&19&71&265&989\\
\hline
7&6&29&139&666&3191\\
\hline
8&7&41&239&1393&8119\\
\hline
9&8&55&377&2584&17711\\
\hline
\end{array}$$

At $N=4$ we obtain $d_k=2k+1$, as for $SO_3$, and at $N\geq5$ the growth is exponential, $d_k\sim cx^k$, with $x+x^{-1}=N-2$. At the 2 special values $N=5$ and $N=9$, the latter corresponding to the color quantum group, we have the following result:

\begin{proposition}
The dimensions of the irreducible representations of $S_5^+$ are the odd Lucas numbers, $d_k=L_{2k+1}$, namely $1,4,11,29,76,\ldots$ The dimensions of the irreducible representations of $S_9^+$, and of its twist $PU_3^+$, are the Fibonacci numbers $d_k=F_{4k+2}$, namely $1,8,55,377,2584,\ldots$
\end{proposition}

\begin{proof}
The Lucas numbers satisfy $L_{n+2}=3L_n-L_{n-2}$, with $L_1=1$, $L_3=4$, which is the above recurrence at $N=5$. The Fibonacci numbers satisfy $F_{n+4}=7F_n-F_{n-4}$, with $F_2=1$, $F_6=8$, which is the above recurrence at $N=9$. The assertion about $PU_3^+$ follows from Theorem 1.4, because a twist preserves the dimensions of the corepresentations. The program {\tt rijkaard.py} checks this as well.
\end{proof}

Thus, the $8$-dimensional representation from section 2 sits inside a Fibonacci sequence, for $S_9^+$ and $PU_3^+$, with the next representation having dimension $55$, and so on, while for $G_3$, which is a twist of $PU_3$, the dimensions are those of $PU_3$, namely $1,8,10,10,27,\ldots$ On the other hand, the twist changes the symmetric structure of the category, and more precisely the flip. Note here that the free pair $S_9^+,PU_3^+$ and the pair $G_3,PU_3$ have different tensor categories, and that the graded flip below concerns the latter pair only, the ordinary flip being not colinear for the free color quantum group. To be more precise, we have:

\begin{proposition}
Under the correspondence $T\to T^\sigma$ from Theorem 1.1, the flip $\Sigma$ of $M_N(\mathbb C)\otimes M_N(\mathbb C)$, which is colinear for $PU_N$, corresponds to the following map of $\mathbb C[\Gamma]\otimes\mathbb C[\Gamma]$, which is colinear for $G_H$:
$$T(e_g\otimes e_h)=d(g,h)e_h\otimes e_g$$
This map satisfies $T^2=1$, so it is a symmetric braiding, and at $N=3$ its coefficients are cube roots of unity.
\end{proposition}

\begin{proof}
We must find $T$ with $\Omega_2^{-1}T\Omega_2=\Sigma$. With $T=\Omega_2\Sigma\Omega_2^{-1}$ we obtain $T(e_g\otimes e_h)=\sigma(g,h)^{-1}\sigma(h,g)e_h\otimes e_g=d(g,h)e_h\otimes e_g$, as claimed. Since $d(g,h)d(h,g)=1$ we have $T^2=1$, and since $d$ is a bicharacter, $T$ satisfies the hexagon relations, so it is a symmetric braiding. At $N=3$ the values of $d$ are the cube roots of unity. Finally, $T$ is colinear for $G_H$ by Theorem 1.1, because $\Sigma$ is colinear for $PU_N$.
\end{proof}

This is quite interesting, historically. Color was introduced by Greenberg in 1964 \cite{gre} as a parastatistics of order $3$ for the quarks, before being reformulated as an $SU_3$ gauge symmetry, and Proposition 7.2 shows that, after untwisting, a $\mathbb Z_3^2$-graded symmetric braiding comes back. We should mention however that this is a statement about the symmetric monoidal structure: the ordinary flip is not colinear for the free color quantum group, the pointwise product on $\mathbb C^9$ still commutes, and no quark exchange operator, Fock representation or parafermion relation is obtained in this way. At the level of Lie algebras, Proposition 7.2 leads to the color Lie algebras of Rittenberg and Wyler \cite{rwy}, also called $\varepsilon$-Lie algebras, studied by Scheunert \cite{sch}. These are $\Gamma$-graded vector spaces $L$ with a bilinear bracket satisfying, for $x\in L_g,y\in L_h,z\in L_k$:
$$[y,x]=-\varepsilon(h,g)[x,y]\quad,\quad\varepsilon(k,g)[x,[y,z]]+\varepsilon(g,h)[y,[z,x]]+\varepsilon(h,k)[z,[x,y]]=0$$

Here $\varepsilon:\Gamma\times\Gamma\to\mathbb T$ is a commutation factor, namely a bicharacter satisfying $\varepsilon(g,h)\varepsilon(h,g)=1$. The basic example is that of the Lie superalgebras, where $\Gamma=\mathbb Z_2$ and $\varepsilon(g,h)=(-1)^{gh}$, and the general theory, including the Poincar\'e--Birkhoff--Witt theorem and the relation with the ordinary Lie superalgebras via cocycle twists, was developed by Scheunert in \cite{sch}. In our situation the grading group is $\Gamma=\widehat{H}\times H$, and the commutation factor is the symplectic form $d$, which takes its values in the $N$-th roots of unity. The point is that $d$ is indeed a commutation factor, because it is a bicharacter satisfying $d(g,h)d(h,g)=1$, and we have:

\begin{proposition}
The group algebra $\mathbb C[\Gamma]$, with the following bracket, is a color Lie algebra, with commutation factor $d$:
$$[e_g,e_h]=e_ge_h-d(g,h)e_he_g=(1-d(g,h))e_{g+h}$$
The Lie algebra $gl_N$ is the twist of this color Lie algebra, with $sl_N$ corresponding to the span of the elements $e_g$ with $g\neq0$, and with the center corresponding to $\mathbb Ce_0$.
\end{proposition}

\begin{proof}
The bracket is $m-m\circ T$, with $T$ being the braiding from Proposition 7.2, so under the correspondence from Theorem 1.1 it becomes $m_\sigma-m_\sigma\Sigma$, which is the commutator of $M_N(\mathbb C)$. Explicitly, in terms of the Weyl operators, we have the following formula:
$$[W_g,W_h]=\sigma(g,h)(1-d(g,h))W_{g+h}$$

The identities of a color Lie algebra are transported from the usual identities for the commutator, as in \cite{sch}. Finally, $d(0,h)=1$, so $e_0$ is central, and $tr(W_g)=0$ for $g\neq0$. We have checked all this at $N=3$ with the program {\tt hagi.py}, including the $\varepsilon$-Jacobi identity for all the $729$ triples of basis elements.
\end{proof}

As for the quarks, Proposition 6.1 shows that a color triplet has no image as a module over the $9$ points. However, $\mathbb C^3$ is the unique irreducible projective representation of $\mathbb Z_3^2$ with cocycle $\sigma$, that is, the unique irreducible module over the twisted group algebra, and the class $[\sigma]\in H^2(\mathbb Z_3^2,\mathbb T)=\mathbb Z_3$ is formally analogous to discrete torsion. This notion was introduced by Vafa \cite{vaf}, see also \cite{vwi}, and Douglas showed in \cite{dou} that it is implemented in D-brane gauge theories by using projective representations of the orbifold group. Making this analogy precise would require a module category, or a field theory, and we will not do this here.

Let us discuss now the gauge fields. A $gl_3$-valued gauge field on a manifold $M$ can be decomposed, locally or on a trivial bundle, in the Weyl basis, which is orthonormal with respect to $<a,b>=tr(ab^*)$, and this decomposition is related to the finest grading of $sl_N$ studied by Patera and Zassenhaus \cite{pz}. We have:

\begin{theorem}
Write a $gl_3$-valued gauge field as $A=\sum_gA_gW_g$, with $A_g$ being $1$-forms on $M$, and with $g\in\mathbb Z_3^2$. The curvature $F=dA+A\wedge A$ has then components
$$F_g=dA_g+\frac{1}{2}\sum_{h+k=g}\sigma(h,k)(1-d(h,k))A_h\wedge A_k$$
and the Yang--Mills density is $\sum_g|F_g|^2$, up to the coupling constant, with respect to $<a,b>=tr(ab^*)$ and to the metric of $M$. The $U_1$ part $A_0$ decouples, and the components of the $8$ gluon fields are labeled by the nonzero elements of $\mathbb F_3^2$, with exactly $48$ ordered pairs of labels having nonzero brackets, namely the pairs of independent labels.
\end{theorem}

\begin{proof}
We have $A\wedge A=\sum_{hk}\sigma(h,k)A_h\wedge A_kW_{h+k}$, and by antisymmetrizing, using $\sigma(k,h)=\sigma(h,k)d(h,k)$, we obtain the formula of $F_g$. The formula of the Yang--Mills density follows from the orthonormality of the Weyl operators. Finally, $d(h,k)=w^{\det(h,k)}$, so $1-d(h,k)=0$ precisely when $h,k$ are dependent. In particular, $F_0=dA_0$, because $h+k=0$ implies that $h,k$ are dependent. For each of the $8$ nonzero labels $h$, there are $6$ nonzero labels $k$ independent from $h$, which gives $48$ pairs.
\end{proof}

With the normalization $J_g=w^{2g_1g_2}W_g$, where $g=(g_1,g_2)$, we obtain the following formula, as checked by {\tt hagi.py}:
$$[J_g,J_h]=2i\sin\left(\frac{2\pi\det(g,h)}{3}\right)J_{g+h}$$

These are the trigonometric structure constants of Fairlie, Fletcher and Zachos \cite{ffz}. Regarding reality, write $A=\sum_gB_gJ_g$, with new coefficients $B_g$. For an $su_3$ connection, in the anti-Hermitian convention, we have $B_0=0$, and since $J_g^*=J_{-g}$, the condition $A^*=-A$ reads $B_{-g}=-\overline{B_g}$, while in the Weyl basis, where $W_g^*=w^{g_1g_2}W_{-g}$, it reads $A_{-g}=-w^{g_1g_2}\overline{A_g}$. Thus the 8 labels pair into 4 complex 1-forms, corresponding to the 8 real color directions. Observe that the labels $g$ index the Weyl basis, which is a Fourier basis with respect to the 9 points, so the points themselves do not label individual gluons. After untwisting, by Proposition 7.3 the structure constants become simply $1-d(h,k)$, the ordinary spacetime 1-forms and the statistics of the quantized fields being unchanged. Making sense of the corresponding untwisted Yang--Mills theory is an interesting question.
Let us compare this with the other twists of the Standard Model geometry which appeared in the literature. In the twisted spectral triples of Connes and Moscovici \cite{cmo}, the algebra is unchanged, and the commutators with the Dirac operator are replaced by twisted commutators $Da-\rho(a)D$, with $\rho$ being an automorphism of the algebra. These were applied to the Standard Model by Devastato and Martinetti \cite{dma}, the breaking of the grand symmetry down to the Standard Model becoming a dynamical process, and were further studied in \cite{fma,lma,msi}, notably in connection with the Lorentzian signature. A different notion, the twisted reality condition, where the real structure is twisted instead, was introduced by Brzezi\'nski, Ciccoli, Dabrowski and Sitarz \cite{bcds}.

These twists are not our cocycle twist, which changes the product of the algebra, and leaves the representation theory of the quantum symmetries unchanged. There is however a formal analogy: for an element $a$ which is an eigenvector of $\rho$, with eigenvalue $\chi$, the twisted commutator is $Da-\chi aD$, which has the same shape as the bracket in Proposition 7.3. In our setting the Dirac operator $D_F$ is color-blind, by Theorem 6.3, so it has degree $0$, and, whenever the transport to the untwisted side exists, its commutators untwist into usual commutators. Thus a link with \cite{dma,lma} would rather involve a Dirac operator, or a real structure, of nonzero degree.

Let us discuss now the weak case. At $N=2$ we have $\Gamma=\mathbb Z_2^2$, the Weyl operators are $1,X,Z,ZX$, with $X,Z$ being Pauli matrices, and $d(g,h)=(-1)^{\det(g,h)}$. Thus, in the weak sector, the graded flip is given by signs, instead of cube roots of unity, and we are led to the $\mathbb Z_2\times\mathbb Z_2$-graded color Lie algebras considered by Rittenberg and Wyler in \cite{rw2}. To be more precise, we have the following result:

\begin{theorem}
The Lie algebra $gl_2$ is the twist of the color Lie algebra $\mathbb C[\mathbb Z_2^2]$, with commutation factor $d(g,h)=(-1)^{\det(g,h)}$, and with bracket given by
$$[e_g,e_h]=2e_{g+h}\quad,\quad\forall g\neq h,\ g,h\neq0$$
and $[e_0,e_h]=[e_g,e_g]=0$. This bracket is symmetric on the nonzero labels, and accordingly, in the Pauli basis, the 3 components of the weak gauge field are labeled by the nonzero elements of $\mathbb F_2^2$, with all the 6 ordered pairs of distinct labels having nonzero brackets.
\end{theorem}

\begin{proof}
This is the case $N=2$ of Proposition 7.2, Proposition 7.3 and Theorem 7.4, with $d(g,h)=-1$ precisely when $g,h$ are distinct and nonzero, so that $1-d(g,h)=2$, and $[e_h,e_g]=-d(h,g)[e_g,e_h]=[e_g,e_h]$. We have checked all this with the program {\tt kluivert.py}, which also checks that the models $P^x$ from section 5 are magic and noncommutative, for random $x\in SU_2$.
\end{proof}

Summarizing, each scalar summand corresponds to a single point, where nothing happens, the electroweak algebra corresponds to $1+1+4$ points, with $S_4^+$ acting on the 4 points and untwisting into $SO_3$, with a graded flip given by signs, and the color sector corresponds to $9$ points, with the free pair $S_9^+,PU_3^+$, and, on restricting to $G_3,PU_3$, with a graded flip given by cube roots of unity. Observe however that electromagnetism, with charge $Q=T_3+Y$, mixes the weak and hypercharge sectors, so it is not modeled by a single scalar summand. The electroweak sector, which is the simplest one, certainly deserves a separate study.

\section{Spectral action and classicality}

We discuss here some general mechanisms forcing quantum symmetries to be classical, which are relevant to the passage from the finite geometry $F$ to the product $M\times F$. The first one is the flip mechanism, used in the proof of Proposition 4.1: if the flip $\Sigma$ intertwines $u^{\otimes2}$, for a faithful corepresentation $u$, then the coefficients of $u$ commute. The second one concerns tensor products:

\begin{proposition}
Let $A_1,A_2$ be unital $C^*$-algebras with coactions $\alpha_i:A_i\to A_i\otimes C(G)$, and assume that a coaction $\alpha$ on the spatial tensor product $A_1\otimes A_2$ satisfies $\alpha(a\otimes1)=\alpha_1(a)_{13}$ and $\alpha(1\otimes b)=\alpha_2(b)_{23}$. Then the coefficients of $\alpha_1$ commute with those of $\alpha_2$. In particular, if both coefficient algebras generate $C(G)$, then $G$ is classical.
\end{proposition}

\begin{proof}
The elements $a\otimes1$ and $1\otimes b$ commute, so their images commute too, and by applying linear forms on the $A_1$ and $A_2$ legs, we obtain the commutation of the coefficients.
\end{proof}

This does not apply to the global actions on $M\times F$, which are trivial on the first leg, but shows that local quantum symmetries on products of independent algebras must be handled with care. Recall also that, by \cite{gsw}, a compact connected smooth manifold has no genuine compact quantum group symmetries, in the sense of faithful smooth actions, so for a product $M\times F$, for an action preserving the manifold factor, the effective action on this factor is classical, although this alone does not classify the actions on the product.

More generally, the passage from $F$ to the product $M\times F$, with algebra $C^\infty(M)\otimes A_F$, Hilbert space $L^2(M,S)\otimes H_F$ and Dirac operator $D_M\otimes1+\gamma_5\otimes D_F$, raises problems at 3 levels. First, the global symmetries are no problem, because the quantum group acts via $id\otimes$coaction, and the cocycle is constant along $M$, so the twisting passes. Second, the local symmetries, or gauge transformations, are maps $M\to G$, and for quantum groups one needs quantum families of maps \cite{sol}, or $C(M)$-linear coactions, as in \cite{bddd}, which is a hard problem. Observe however that $H^2(\mathbb Z_3^2,\mathbb T)=\mathbb Z_3$ is discrete, so on a connected manifold $M$ a locally defined twist class is forced to be constant, but this does not produce a global Weyl frame on a nontrivial bundle, so twisting fiberwise is fine only locally, or on trivial bundles. Third, there is the question of the fermions, discussed in section 6, which is where the real obstruction lies. Observe also that, at the level of the gauge fields, the untwisting can be performed locally, or on trivial bundles, as in Theorem 7.4, with the color Lie algebra from Proposition 7.3 replacing $gl_3$.

The third mechanism concerns the spectral action \cite{cc2}. Consider a real spectral triple $(A,H,D,J)$, and a finite quantum group $Q$, with $C(Q)$ acting on $H_Q=L^2(C(Q),h)$ by left multiplication, and with $j(a)=a^*$, and set $K=J\otimes j$. Consider a coaction implemented by a unitary $U$ on $H\otimes H_Q$, which commutes with $D\otimes1$. There are 2 natural reality conditions for $U$, namely the full commutation $KU=UK$, and the following coefficientwise condition, used in \cite[Definition 2.12]{bdd}:
$$KU(\xi\otimes1)=U(J\xi\otimes1)\quad,\quad\forall\xi\in H$$

Given a classical self-adjoint $1$-form $A_1=\sum a_i[D,b_i]$, we have the fluctuated operator $D_{A_1}=D+A_1+\varepsilon'JA_1J^{-1}$, and its transformed version $D'=D\otimes1+A_1'+\varepsilon'KA_1'K^{-1}$, with $A_1'=U(A_1\otimes1)U^*$. The natural invariance requirement for the bosonic spectral action, with cutoff function $f$ and scale $\Lambda$, is then the following equality of operators on $H_Q$:
$$(Tr_H\otimes id)f(D'/\Lambda)=Tr_Hf(D_{A_1}/\Lambda)1$$

Under full commutation we have $D'=U(D_{A_1}\otimes1)U^*$, and the invariance follows by functional calculus, and the Kac type partial trace identity from the proof of Proposition 3.2. Under the coefficientwise condition this fails, because $KL_aK^{-1}$ involves a right multiplication by $a^*$, and we have the following result, which shows that the invariance then detects all the nonclassical finite quantum groups:

\begin{theorem}
For a finite quantum group $Q$, the following are equivalent:
\begin{enumerate}
\item $C(Q)$ is commutative.

\item The above bosonic invariance holds for all real product geometries, coefficientwise real actions, classical $1$-forms and cutoffs.

\item The bosonic invariance holds for the regular coaction of $Q$, with $1$-forms built from the counit block projection, on a fixed $2$-point external triple, or alternatively on the flat torus $\mathbb T^4$, each of these 2 geometries being sufficient.
\end{enumerate}
\end{theorem}

\begin{proof}
Write $C(Q)=\oplus_\beta M_{n_\beta}(\mathbb C)$, with central projections $z_\beta$, and $d=\sum_\beta n_\beta^2$. Let $e$ be the counit block projection, $\Pi=\Delta(e)$, and $P=L_\Pi$, $R=R_\Pi$, acting on $H_Q\otimes H_Q$. For the regular coaction we have $U(L_e\otimes1)U^*=P$ and $KPK^{-1}=R$. In the block $\alpha\otimes\beta$, the element $\Pi$ projects onto the invariant vectors of $\pi_\alpha\otimes\pi_\beta$, which by Schur's lemma form the line spanned by a unit maximally entangled vector when $\beta$ is the dual of $\alpha$, and are zero otherwise. In the Hilbert--Schmidt realization the map $X\to\Pi X\Pi$ is then a rank one projection, whose vector has $n_\beta^2$ equal Schmidt coefficients $1/n_\beta$, and with $E$ being the ordinary partial trace over the first leg, we obtain:
$$E(PR)=\sum_\beta n_\beta^{-2}L_{z_\beta}\quad,\quad E(P)=E(R)=1\quad,\quad E(1)=d$$

For the 2-point witness, represent $\mathbb C^2$ on $\mathbb C^3$, with basis $0,L,R$ carrying the bimodule labels $(1,1),(2,1),(1,2)$, let $J$ exchange $L,R$ and conjugate the coefficients, and set $B=|0\rangle\langle L|+|L\rangle\langle0|$, $C=JBJ^{-1}$, $D=B+C$ and $\gamma=diag(1,-1,-1)$. This is a real even spectral triple satisfying the first order condition, and with $p=diag(0,1,0)$ we have $B=(1-2p)[D,p]$. Take its product with $(C(Q),H_Q,0,j)$, with the regular coaction, which satisfies the coefficientwise condition, and the $1$-form $-B\otimes L_e$. Since $e$ is central, the fluctuated operator is $(B+C)\otimes(1-L_e)$, whose kernel has dimension $d+2$. After transformation we obtain $B\otimes(1-P)+C\otimes(1-R)$, and on a joint eigenspace of $P,R$ the kernel has dimension $1$, unless $P=R=1$ there, where it has dimension $3$. With a cutoff selecting the kernel, the transformed action is therefore $d+2E(PR)$, which equals $d+2/n_\beta^2$ on the block $\beta$, so the invariance fails as soon as some $n_\beta>1$. The same works with the heat cutoff $f(D)=e^{-sD^2}$, at any scale $s>0$, the transformed action minus the original one being $2(1-e^{-s})^2(n_\beta^{-2}-1)$ on the block $\beta$.

On the $2\pi$-periodic flat torus $\mathbb T^4$, with periodic spin structure and $D=-i\sum\gamma^\mu\partial_\mu$, use instead the $1$-form $t\gamma^1\otimes L_e$, with $0<t<1/2$, coming from $u^*[D,u]=\gamma^1$ with $u=e^{ix_1}$. The original shifts cancel, while after transformation the first momentum is shifted by $t(P-R)$, and only the zero momentum with $P=R$ contributes to the kernel, which gives $4(d-2+2E(PR))$ instead of $4d$. Conversely, in the classical case the action is fiberwise given by unitaries commuting with $D$ and $J$, and the invariance holds. We have checked all this with the programs {\tt beckham.py} and {\tt zlatan.py}, the latter building the transformed operators, and computing the heat cutoff version, for the duals of $S_3$ and $Q_8$.
\end{proof}

The torus witness shows in fact more than the failure of invariance. Indeed, the 1-forms $A_1=0$ and $A_1=t\gamma^1\otimes L_e$ give the same fluctuated operator $D_{A_1}$, but different transformed operators, so the prescribed transformation of the 1-forms does not descend to the fluctuated Dirac operators.

For the Standard Model geometry itself, a natural question is whether the invariance can be restored by using a second, $2$-sided implementing unitary $W$, commuting with $K$ and with the Dirac operator. Observe that if $U$ has coefficients acting by left multiplications, and $KU=UK$, then these coefficients lie in $L(C(Q))\cap R(C(Q))$, which is the center of $C(Q)$, because $L_x=R_y$ gives $x=y$ by evaluating at $1$, and then $xz=zx$ for any $z$. Thus such a $W$ must be of a different type. For the Standard Model, such a $W$ exists, the point being that the weak algebra acts on the particle doublets, while the color algebra acts on the antiquarks, so that the weak phase and the color action can be placed on opposite legs of $L^2(Q)$. To be more precise, we have the following result:

\begin{proposition}
For the quantum groups $Q$ from Theorem 2.4, acting on $F$, there is a unitary $W$ on $H_F\otimes L^2(Q)$, spatial in the Haar representation, which commutes with $D_F\otimes1$, $\gamma_F\otimes1$ and $K$, implements the effective coaction on the represented complex algebra $B_F$, and satisfies $(Tr_{H_F}\otimes id)(W(X\otimes1)W^*)=Tr(X)1$. Thus, for classical 1-forms, the bosonic invariance holds for $F$, and for its products with spin manifolds, with the global action, trivial on the manifold, and with cutoffs giving trace-class operators.
\end{proposition}

\begin{proof}
Let $\mathcal M=L^\infty(Q)$, which is a finite von Neumann algebra, $Q$ being of Kac type, and write $L_a,R_a$ for the left and right multiplications on $L^2(Q)$, so that $jL_aj=R_{a^*}$. The color coaction is a trace preserving morphism $\beta:M_3(\mathbb C)\to M_3(\mathbb C)\otimes\mathcal M$, and the projections $\beta(e_{11})$ and $e_{11}\otimes1$ both have center-valued trace $1/3$, so there is a partial isometry $s$ with $s^*s=e_{11}\otimes1$ and $ss^*=\beta(e_{11})$. Then $v=\sum_i\beta(e_{i1})s(e_{1i}\otimes1)$ is a unitary with $\beta(b)=v(b\otimes1)v^*$, and trace preservation gives $\sum_iv_{ij}v_{ik}^*=\delta_{jk}1$. With $z$ being the generator of the weak circle, set $h_\nu=h_u=1$ and $h_e=h_d=z^*$, and let $W$ act on the particles of species $s$ by $L_{h_s}V_R$, and on the antiparticles of species $s$ by $R_{h_s^*}V_L$, where $V_L=(L_{v_{ij}})$ and $V_R=(R_{v_{ij}^*})$ act on the color indices of the quarks, and are trivial on the leptons. Then $W$ is unitary, and commutes with $K$, because $jL_aj=R_{a^*}$. On the weak doublets $V_R$ cancels, and we obtain $e_{12}\to e_{12}\otimes z$, while on the antiquarks the right phases commute with the left color action, and we obtain $\beta$, so $W$ implements the coaction. Since $W$ acts in the same way on both ends of each mass arrow, and trivially on the generations and on the Majorana arrow, it commutes with $D_F$ and $\gamma_F$. The partial trace formula follows from $\sum_iw_{ij}w_{ik}^*=\delta_{jk}1$ for the color blocks $w$ of $W$, which comes from the above relation for $v$, and from $R_aR_b=R_{ba}$. Finally, $W$ commutes with $K$ and with the Dirac operator, so $D'=W(D_{A_1}\otimes1)W^*$, and the partial trace formula gives the invariance. We have checked all this with the program {\tt vanbasten.py}, in a model with $\mathcal M=\mathbb C[S_3]$, 3 generations, generation mixing and a Majorana term.
\end{proof}

Observe that $W$ is not claimed to be a corepresentation, and that the compatibility with the real form of $A_F$, which selects $QISO^+_{\mathbb R}$ in Theorem 2.4, is an additional requirement, not used here.

This is in agreement with \cite[section 4]{bdd}, where the invariance of the spectral action under the quantum isometries of $F$ is established by using the conjugacy identity \cite[(4.4)]{bdd}, which follows from the full commutation with $K$, but not from the coefficientwise condition alone, as the above regular examples show. Thus such statements depend on the precise prescription for transforming the fluctuations, on the reality condition imposed on the implementer, and on the geometry itself, the obstruction from Theorem 8.2 coming from the regular coactions.

\section{Mixing matrices and Hadamard}

The Standard Model has 2 mixing matrices, namely the CKM matrix for the quarks, relating the mass eigenstates of the up and down quarks, and the PMNS matrix for the leptons, relating the mass and flavor eigenstates of the neutrinos. Both are unitaries $C\in U_3$, defined up to phases on rows and columns, at least as far as the oscillations are concerned, 2 Majorana phases of the PMNS matrix being physical in the Majorana case, but invisible in the oscillations, and in the Chamseddine--Connes framework they appear inside the Dirac operator $D_F$, via the Yukawa couplings, see \cite{ccm,con}.

In \cite{co2}, Connes interpreted such a matrix as the relative position of 2 maximal abelian subalgebras of $M_n(\mathbb C)$, namely the diagonal algebra $D$ and its conjugate $UDU^*$, and extended this into a unitary invariant of Riemannian manifolds, which together with the spectrum of the Dirac operator is a complete invariant, see also \cite{co3}. A key example in \cite{co2}, in the section on the CKM matrix and the Fourier transform, is the Fourier matrix $F_3=(w^{ij})$, with $w=e^{2\pi i/3}$, which is said to fall in the class of commuting squares of Popa \cite{pop}.

Let us recall that 2 subalgebras $D_1,D_2\subset M_n(\mathbb C)$, together with a common subalgebra $D_0\subset D_1\cap D_2$, form a commuting square with respect to the trace when $E_{D_1}E_{D_2}=E_{D_0}$, the maps $E$ being the conditional expectations, see \cite{pop}. In the case of maximal abelian subalgebras, with $D_0=\mathbb C1$, which corresponds to the setting of Connes, who works with labeled minimal idempotents, we have the following well-known result, see \cite{haa,pop,tzy}:

\begin{proposition}
For a unitary $U\in U_n$, we have $E_DE_{UDU^*}=E_{\mathbb C1}$ if and only if $|U_{ij}|^2=\frac{1}{n}$ for any $i,j$, that is, if and only if $\sqrt{n}U$ is a complex Hadamard matrix. In this case we have as well $E_{UDU^*}E_D=E_{\mathbb C1}$, and $D\cap UDU^*=\mathbb C1$.
\end{proposition}

\begin{proof}
Let $e_i\in D$ be the minimal projections, and $f_j=Ue_jU^*$. Since $E_{\mathbb C1}=E_{\mathbb C1}E_{UDU^*}$, the equality to be proved holds if and only if it holds on each $f_j$. But we have $E_D(f_j)=\sum_in\,tr(e_if_j)e_i=\sum_i|U_{ij}|^2e_i$, while $E_{\mathbb C1}(f_j)=\frac{1}{n}1$, which gives the result. The other assertions follow by taking adjoints, and from $E_D(x)=x$ for $x\in D\cap UDU^*$.
\end{proof}

When formulating instead the commuting square condition over $D\cap UDU^*$, the identity matrix gives a trivial example. In general, the bipartite graph having as edges the pairs $(i,j)$ with $U_{ij}\neq0$ decomposes $U$, after permuting the rows and columns, into square blocks, the intersection $D\cap UDU^*$ being scalar on each block, and the commuting square condition holds precisely when each block is, up to normalization, a complex Hadamard matrix of its own size.

Thus, the commuting square case of the Connes invariant is precisely the complex Hadamard case. Now recall that any complex Hadamard matrix $H\in M_N(\mathbb C)$ produces a quantum permutation group $G\subset S_N^+$, as follows, see \cite{bhm}: the vectors $\xi_{ij}=H_i/H_j$, obtained by dividing the rows of $H$ entrywise, form a magic basis of $\mathbb C^N$, the magic matrix $P_{ij}=Proj(\xi_{ij})$ defines a representation $\pi:C(S_N^+)\to M_N(\mathbb C)$, and $G$ is given by the Hopf image of $\pi$, that is, $C(G)$ is the smallest Hopf algebra quotient of $C(S_N^+)$ through which $\pi$ factors.

Regarding now the physics, the quark mixing matrix CKM is close to the identity, which is the opposite of the Hadamard regime. However, the lepton mixing matrix PMNS is much closer to the Hadamard regime. In fact, the Fourier matrix $F_3/\sqrt{3}$ itself was proposed in 1978 by Cabibbo and by Wolfenstein as a lepton mixing matrix, under the name trimaximal mixing \cite{cab,wol}. The more recent tribimaximal ansatz of Harrison, Perkins and Scott \cite{hps} is as follows, up to signs and phases:
$$U_{TB}=\begin{pmatrix}\sqrt{2/3}&1/\sqrt{3}&0\\ -1/\sqrt{6}&1/\sqrt{3}&1/\sqrt{2}\\ -1/\sqrt{6}&1/\sqrt{3}&-1/\sqrt{2}\end{pmatrix}$$

The middle column here is flat, as for a Hadamard matrix. The matrix $F_3/\sqrt{3}$, which predicts $|U_{e3}|^2=1/3$, was ruled out by the data long ago, and the measurement of a nonzero reactor angle $\theta_{13}$ by Daya Bay and others \cite{dby} ruled out $U_{TB}$ too. As for the trimaximal scheme TM2, keeping the flat middle column of $U_{TB}$, this predicts $\sin^2\theta_{12}=1/(3\cos^2\theta_{13})\simeq0.341$, with $\sin^2\theta_{13}\simeq0.022$, and the first precision measurement of $\theta_{12}$ by JUNO \cite{jun} gives a value close to $0.31$, with an uncertainty below $0.01$, which puts exact TM2 under strong tension, leaving only approximate textures, with explicit deviations. Thus the lepton mixing matrix is much closer to the Hadamard regime than the quark mixing matrix, and we will make this more precise at the end of this section.

Most models producing tribimaximal or related mixing use a finite flavor group, acting on the 3 generations via a 3-dimensional representation. The classic examples are the tetrahedral group $A_4$, introduced by Ma and Rajasekaran \cite{mra} and used by Altarelli and Feruglio \cite{afe} for tribimaximal mixing, the octahedral group $S_4$, and the icosahedral group $A_5$, which produces the golden ratio mixing \cite{est}. These, together with the cyclic and dihedral groups, are the finite subgroups of $SO_3$, acting on $\mathbb R^3$.

Now recall from section 3 that $SO_3$ is the twist of the weak quantum group $S_4^+$. The quantum subgroups of $S_4^+$ were classified in \cite{bbi}, by using the isomorphism $S_4^+\simeq SO_3^{-1}$, and include, besides $S_4$ and its subgroups, the twists $S_4^\tau$ and $A_5^\tau$ of the octahedral and icosahedral groups, which are not classical, and are subject to a McKay type correspondence. Thus, the classical flavor groups have quantum permutation counterparts, acting on 4 points, and the integration theory over $S_4^+$, developed in \cite{bcp}, could be useful here.

There is also a link with the fermions, discussed in section 6. The group $\Delta(27)$ acts on the 3 generations in many flavor models \cite{bgg,hls}, and its 3-dimensional representations are generated by the clock and shift matrices, the Fourier matrix $F_3$ interchanging their eigenbases, which are therefore mutually unbiased. Thus, the same matrix $F_3$ appears in the color twist from section 1, in the Connes example from \cite{co2}, and in the trimaximal mixing from \cite{cab,wol}. Summarizing, we have the following result:

\begin{theorem}
The Connes invariant of a pair of maximal abelian subalgebras $D,UDU^*\subset M_n(\mathbb C)$ corresponds to a commuting square over $\mathbb C1$ precisely when $\sqrt{n}U$ is a complex Hadamard matrix, and in this case it produces a quantum permutation group $G\subset S_n^+$. For the trimaximal matrix $F_3/\sqrt{3}$ this quantum permutation group is the group $\mathbb Z_3$.
\end{theorem}

\begin{proof}
The first assertion is Proposition 9.1. Regarding the second assertion, for a Fourier matrix $F_G$ we have $\xi_{ij}=\xi_{i-j,0}$, so the magic matrix is circulant, and its coefficients generate the algebra $C(G)$, as explained in \cite{bhm}. At $G=\mathbb Z_3$ we obtain the result.
\end{proof}

Let us discuss now the color case in more detail. At $N=3$, the 8 nonzero labels of the Weyl operators $W_g$ split into 4 pairs $\{g,2g\}$, corresponding to the 4 lines through $0$ in $\mathbb F_3^2$, and we have:

\begin{proposition}
For each of the 4 lines $L$ through $0$ in $\mathbb F_3^2$, the operators $W_g$ with $g\in L$ have a common eigenbasis $B_L$ of $\mathbb C^3$, unique up to phases and permutations. These 4 bases are mutually unbiased, and the transition matrix between any 2 of them, multiplied by $\sqrt{3}$, is a complex Hadamard matrix, equivalent to $F_3$.
\end{proposition}

\begin{proof}
For $g\neq0$ we have $W_{2g}\sim W_g^2$, and $W_g$ has the 3 cube roots of unity as eigenvalues, up to a common phase, because $tr(W_g)=tr(W_g^2)=0$. This gives the bases $B_L$. If $g,h$ are independent, $W_hW_g\sim W_gW_h$ with a nontrivial cube root of unity as scalar, so $W_h$ permutes cyclically the eigenvectors $e_i$ of $W_g$, up to phases. Thus for an eigenvector $f$ of $W_h$ the quantities $|<e_i,f>|$ do not depend on $i$, and so are all equal to $1/\sqrt{3}$. Finally, any $3\times3$ complex Hadamard matrix is equivalent to $F_3$, see \cite{haa,tzy}.
\end{proof}

These are the mutually unbiased bases of Ivanovi\'c and Wootters--Fields \cite{ivi,wfi}. Explicitly, consider the Fourier matrix and the phase matrix, given by:
$$F_3=\begin{pmatrix}1&1&1\\ 1&w&w^2\\ 1&w^2&w\end{pmatrix}\quad,\quad D=\begin{pmatrix}1&0&0\\ 0&1&0\\ 0&0&w\end{pmatrix}$$

Then, up to phases and to the order of the vectors, the 4 bases $B_L$ with $L$ spanned by $(1,0),(0,1),(1,2),(1,1)$ are given by the columns of the following unitary matrices:
$$1_3\quad,\quad\frac{F_3}{\sqrt{3}}\quad,\quad\frac{DF_3}{\sqrt{3}}\quad,\quad\frac{D^2F_3}{\sqrt{3}}$$

Thus the transition matrices between the standard basis and the other 3 bases are the trimaximal matrix $F_3/\sqrt{3}$ and its versions twisted by the phases of $D$, and the transition matrices between 2 nonstandard bases are of the form $F_3^*D^kF_3/3$ with $k=1,2$, which are complex Hadamard matrices too, after multiplication by $\sqrt{3}$, as predicted by Proposition 9.3. Observe also that the 9 points of $\mathbb F_3^2$, together with the $12$ affine lines, form the Hesse configuration of the 9 inflection points of a plane cubic, the Heisenberg group of order $27$ acting by clock and shift on the Hesse pencil $x^3+y^3+z^3+txyz=0$, and the group of order 216 from Proposition 4.4 being the Hessian group, see \cite{ado}. In connection with the mixing matrices, observe that the bases $B_L$ are well-defined only as sets of rays, the ordered eigenbases being unique up to phases and permutations, so the natural objects here are the 4 sets of rays, corresponding to the 4 parallel classes of lines of $\mathbb F_3^2$. The point is that the Clifford group permutes these 4 sets of rays, and that the flavor groups $A_4$ and $S_4$ appear in this way, as follows:

\begin{theorem}
The projective Clifford group $C_3\subset PU_3$, which is the twisted counterpart of the classical version of $G_3$, of order $216$, permutes the 4 mutually unbiased bases $B_L$, and its action factors as
$$C_3\to C_3/K\simeq SL_2(3)\to PSL_2(3)\simeq A_4$$
via the action of $SL_2(3)$ on the 4 lines through $0$. When adding complex conjugation, which corresponds to orientation reversal, we obtain a group of order $432$, acting on the 4 bases via $GL_2(3)\to PGL_2(3)\simeq S_4$.
\end{theorem}

\begin{proof}
An element $U\in C_3$ satisfies $UW_gU^*\sim W_{Ag}$, where $A\in SL_2(3)$ is its image in the exact sequence from section 5, so $U$ maps $B_L$ to $B_{AL}$. The Weyl group $K$ acts trivially, as does $-1$, and the kernel of the action of $SL_2(3)$ on the 4 points of the projective line over $\mathbb F_3$ is $\{\pm1\}$, with image $PSL_2(3)\simeq A_4$. Complex conjugation maps $W_{(i,j)}$ to $W_{(-i,j)}$, which corresponds to a matrix of determinant $-1$, and we obtain in this way $GL_2(3)$, acting via $PGL_2(3)\simeq S_4$. We have checked all this with the program {\tt cantona.py}, which constructs the 4 bases, verifies that the Fourier matrix and the phase gate $diag(1,1,w)$ normalize the Weyl group and induce generators of $SL_2(3)$, and computes the images of $SL_2(3)$ and $GL_2(3)$ in the symmetric group on the 4 bases, with element orders matching those of $A_4$ and $S_4$.
\end{proof}

This is quite interesting, physically. In the flavor models, the mixing matrix appears as the transition matrix between the eigenbases of the residual symmetries preserved by the charged leptons and by the neutrinos, and it was shown by Lam \cite{lam} that tribimaximal mixing, for all values of the Yukawa couplings, requires the horizontal group $S_4$, or a group containing it, while $A_4$ naturally leads to trimaximal mixing, with the qualification, explained by Grimus, Lavoura and Ludl \cite{gll}, that the symmetries of the effective mass terms need not be symmetries of the underlying Lagrangian. In the above picture, if the residual symmetries are the cyclic subgroups of the Weyl group, the trimaximal matrix $F_3/\sqrt{3}$ is the transition matrix between the eigenbases of clock and shift, and $A_4$ and $S_4$ are the permutation groups of the 4 mutually unbiased bases, induced by the classical versions of the quantum groups from section 4. Let us stress however that this provides no mechanism fixing the eigenbases of the physical Yukawa couplings: the CKM and PMNS matrices can vary continuously, with the color construction being unchanged, so the above is only a conditional statement, about mass matrices having specified residual Weyl symmetries.

At the quantum level, the 4 bases $B_L$ correspond to the 4 parallel classes of lines of $\mathbb F_3^2$, and a natural question is whether the color quantum group $G_3$ acts on these 4 classes, which would give a link with the weak quantum group $S_4^+$. The answer is negative, and in fact we have the following result:

\begin{proposition}
The quantum group $G_3$ has no nontrivial coaction on a set of $k\leq8$ points. In particular, the action of its classical version on the 4 mutually unbiased bases, via $A_4$, does not extend to $G_3$.
\end{proposition}

\begin{proof}
A coaction of $G_3$ on $k$ points gives a $k$-dimensional magic corepresentation, having the constant vector as fixed vector. Since $G_3$ is a twist of $PU_3$, its irreducible corepresentations have the same dimensions as those of $PU_3$, namely the dimensions $(a+1)(b+1)(a+b+2)/2$ of the irreducible representations of $SU_3$ having highest weight $(a,b)$ with $a=b$ modulo 3. The smallest nontrivial such dimension is $8$, attained by the adjoint representation, so for $k\leq8$ our corepresentation is a sum of trivial ones, and so the coaction is trivial. The program {\tt gullit.py} lists these dimensions, namely $1,8,10,27,28,35,\ldots$
\end{proof}

Thus, the passage to the quantum symmetry loses here a classical quotient action, and the bound $k\leq8$ is sharp, as shown by the action on the 9 points. As a consequence, a link between the color and weak sectors, via mixing, must be searched elsewhere, for instance at the level of the classical versions.

Finally, as a classical baseline for the deviation from the Hadamard regime, for a mixing matrix $U\in U_n$ we can use the following rephasing invariant quantity, which vanishes precisely in the Hadamard case, and equals $n-1$ for a permutation matrix:
$$\delta(U)=\sum_{ij}\left(|U_{ij}|^2-\frac{1}{n}\right)^2=\sum_{ij}|U_{ij}|^4-1$$

Observe that $\delta(U)=||E_DE_{UDU^*}-E_{\mathbb C1}||^2$, with the Hilbert--Schmidt norm of linear maps on $M_n(\mathbb C)$, so $\delta$ measures the failure of the commuting square condition from Proposition 9.1. At $n=3$ we have $\delta=2$ for the identity, $\delta=1/3$ for $U_{TB}$, and $\delta=0$ for the transition matrices between the mutually unbiased bases, as checked with the program {\tt overmars.py}, which also checks Proposition 9.1, the above Hilbert--Schmidt formula, and the unitarity of $U_{TB}$. For the CKM matrix, which is close to the identity, $\delta$ is close to its maximal value $2$, while for the PMNS matrix, which is close to $U_{TB}$, it is close to $1/3$, a precise comparison requiring the experimental values, with their uncertainties. Observe that $\delta$ is unchanged by complex conjugation, so it loses the sign of the Jarlskog invariant $J$, and that it compresses even the information carried by the moduli, which determine $|J|$. Summarizing, the question of understanding the mixing matrices in terms of quantum permutations is natural, with the Hadamard case being well understood, and with the deviation from the Hadamard regime remaining to be measured by a suitable quantum group invariant, refining $\delta$.

\section{A look into Yang-Mills theory}

The Yang--Mills theory with gauge quantum group $S_N^+$, or its twists, should normally come only after the untwisting of the whole Standard Model, discussed in sections 3 and 7, and after the passage from $F$ to $M\times F$. Indeed, there is no point in formulating a gauge theory for objects which are still twisted. We gather here, as a program, a few observations which make sense independently of this.

In 2 dimensions, Yang--Mills theory on a closed surface $\Sigma_g$ of genus $g$ and area $A$ can be solved exactly, and the Migdal--Witten formula \cite{mig,wit} gives the partition function as a sum over the irreducible representations, with $c$ being the Casimir:
$$Z(\Sigma_g,A)=\sum_\lambda(\dim\lambda)^{2-2g}e^{-A\,c(\lambda)}$$

As explained by Witten in \cite{wi2}, this sum encodes the intersection pairings on the moduli spaces of flat connections, and at $A=0$ it becomes the Witten zeta function $\sum_\lambda(\dim\lambda)^{-s}$, evaluated at $s=2g-2$. Since this formula only uses the representation theory, one can formally write it for $S_N^+$, and the result is invariant under twisting. With the dimensions $d_k$ from section 7, and with the Casimir $c(k)=k(k+1)$ of the $SO_3$ case, we are led to the following formal quantities:
$$Z_N(\Sigma_g,A)=\sum_kd_k^{2-2g}e^{-Ak(k+1)}\quad,\quad\zeta_N(s)=\sum_kd_k^{-s}$$

The program {\tt rijkaard.py} computes these quantities numerically, and here are the values of $\zeta_N$ at $s=2,4$, corresponding to the genus $g=2,3$, with $4$ decimals:
$$\begin{array}{|c|c|c|c|c|c|c|}
\hline
N&4&5&6&7&8&9\\
\hline
\zeta_N(2)&1.2337&1.0722&1.0430&1.0290&1.0210&1.0160\\
\hline
\zeta_N(4)&1.0147&1.0040&1.0016&1.0008&1.0004&1.0002\\
\hline
\end{array}$$

At $N=4$ these are the values $\pi^2/8$ and $\pi^4/96$, coming from the $SO_3$ theory, and at $N\geq5$ the values approach $1$ quickly, the main correction being $(N-1)^{-s}$, coming from the representation $r_1$. In general, we have the following result:

\begin{theorem}
The formal $S_N^+$ partition function $Z_N(\Sigma_g,A)$, with $N\geq4$, is well-defined for any $A>0$ and any genus $g$, and as $N\to\infty$, with $A>0$ fixed, it diverges for $g=0$, is independent of $N$ for $g=1$, and tends to $1$ for $g\geq2$. At $A=0$ the series $\zeta_N(2g-2)$ converges for $g\geq2$, with exponential speed when $N\geq5$, at $N=4$ we have $\zeta_4(s)=(1-2^{-s})\zeta(s)$ for $s>1$, as for $SO_3$, and $\zeta_N(s)\to1$ as $N\to\infty$, for any $s>0$.
\end{theorem}

\begin{proof}
For $A>0$ the Gaussian factor $e^{-Ak(k+1)}$ dominates the exponential growth of $d_k^{2-2g}$. Since $d_{k-1}\leq d_k$, the recurrence gives $d_{k+1}\geq(N-3)d_k$, and so $d_k\geq(N-1)(N-3)^{k-1}$ for $k\geq1$. Thus for $N\geq5$ and $s>0$ the series $\sum_{k\geq1}d_k^{-s}$ converges exponentially fast, and is bounded by $(N-1)^{-s}\sum_{k\geq0}(N-3)^{-sk}$, which tends to $0$ as $N\to\infty$. At $N=4$ we have $\sum_k(2k+1)^{-s}=(1-2^{-s})\zeta(s)$. Finally, for $g=0$ we have $Z_N\geq(N-1)^2e^{-2A}$, for $g=1$ we have $Z_N=\sum_ke^{-Ak(k+1)}$, and for $g\geq2$ we have $1\leq Z_N\leq\zeta_N(2g-2)$, which tends to $1$.
\end{proof}

This is to be compared with the large $N$ results for $U_N$ and $SU_N$, where the limit of the partition function was computed by Lemoine \cite{lem}, in terms of the Witten zeta function, the Jacobi theta function and the Euler function, and where the Wilson loops were shown by Dahlqvist and Lemoine \cite{dl1,dl2} to converge to the planar master field. Understanding the analogues of these results for $S_N^+$, where the large $N$ limit of the partition function is trivial in genus $g\geq2$, but where the Wilson loops should remain interesting, is a natural question.

On the lattice, in the sense of Wilson \cite{wil}, one can try to use magic unitaries on the links. By Proposition 8.1, commuting link algebras, each generating the whole coefficient algebra, would force the quantum group to be classical, which suggests using free products of link algebras. The Haar state then integrates the link variables via the Weingarten formula, with noncrossing partitions \cite{bcp}. For $S_4^+$ the full Pauli model is available, while at 9 points the Weyl models from section 5 factor through $G_3$, whose Haar measure differs from that of $S_9^+$, already at the level of the third moment of the magic character, which is $6$ for $G_3$, and $5$ for $S_9^+$, so that integrating over unrestricted $S_9^+$ links requires its own Weingarten formula. However, a lattice theory would also need edge reversal, vertex gauge transformations, an action or a measure, and a prescription for the Wilson loops. For a general framework for lattice gauge fields in noncommutative geometry, see the gauge networks of Marcolli and van Suijlekom \cite{mvs}.

Finally, regarding large $N$, recall that in the limit of 't Hooft \cite{tho} the planar diagrams dominate, and that the corresponding $1/N$ expansions of matrix models are now well understood, see for instance \cite{bgu}. Planar diagrams are noncrossing, and so are governed by free probability, and in 2 dimensions, for the master field on the plane, the holonomy along a simple loop of area $t$ is distributed as the free unitary Brownian motion of Biane \cite{bia} at time $t$, see \cite{ase,lev}. On the other hand, the representation theory of $S_{N^2}^+$ is governed by the noncrossing partitions at any $N$, and the Weingarten formula for $S_{N^2}^+$ computes integrals as sums over noncrossing partitions with exact $N$-dependence, while the corresponding formulae for $SU_N$ only become noncrossing as $N\to\infty$. However, the noncrossing partitions appearing in the intertwiner spaces and the planar ribbon graphs of the gauge theory expansions count different objects, as already shown at finite rank by the different cubic multiplicities in Proposition 2.3, and relating them, if possible, is part of the program.

Let us also record what happens in the 2 cases of interest for the Standard Model. In the weak case $N=4$ the dimensions are $d_k=2k+1$, and the $S_4^+$ partition function coincides with the partition function of the $SO_3$ theory, as it should, because $S_4^+$ is a twist of $SO_3$. In the color case $N=9$ the dimensions are the Fibonacci numbers $F_{4k+2}$, from Proposition 7.1, this concerning $S_9^+$ and its twist $PU_3^+$, while the color group $PU_3$ of the real case, and its twist $G_3$, have the representation theory of $PU_3$. The zeta function becomes then a Fibonacci zeta function, as follows:
$$\zeta_9(s)=\sum_{k\geq0}\frac{1}{F_{4k+2}^s}=1+\frac{1}{8^s}+\frac{1}{55^s}+\frac{1}{377^s}+\ldots$$

Since $F_{4k+2}\sim\varphi^{4k+2}/\sqrt{5}$, with $\varphi$ being the golden ratio, this series converges for any $s>0$, and is dominated by its first terms. In particular, in genus $g\geq2$ the color partition function at $A=0$ is very close to $1$, the correction coming from the $8$-dimensional representation, of order $8^{2-2g}$. This is in contrast with the $SU_3$ theory, where the dimensions grow only polynomially, the Witten zeta function being of number theoretic nature.

Let us point out, however, that the above Casimir weights, coming from the $SO_3$ case, do not come from a positive heat kernel on $S_N^+$ at $N\geq5$, and so cannot be used as such for the Yang--Mills measure. To be more precise, we have the following result:

\begin{theorem}
For $N\geq5$, the linear functionals $\varphi_t$ on the dense Hopf $*$-algebra of $C(S_N^+)$ given by $\varphi_t(u^{(k)}_{ij})=\delta_{ij}e^{-tk(k+1)}$, where $u^{(k)}$ are the irreducible corepresentations, are not positive, for $t>0$ small.
\end{theorem}

\begin{proof}
Let $\chi$ be the character of $r_1$, which is self-adjoint, with $\chi^2=1+\chi+\chi_2$, and set $d=N-1$, so that $\dim r_2=d^2-d-1$. If $\varphi_t$ is a state, the variance of $\chi$ must be $\geq0$, and this variance is given by:
$$V(t)=1+de^{-2t}+(d^2-d-1)e^{-6t}-d^2e^{-4t}$$

But we have $V(0)=0$ and $V'(0)=-2d-6(d^2-d-1)+4d^2=-2N(N-4)$, so $V(t)<0$ for $t>0$ small, when $N\geq5$. We have checked this with the program {\tt stoichkov.py}, which also computes the large $N$ behavior from Theorem 10.1.
\end{proof}

The problem is the quadratic growth of the Casimir weights. In fact, the positive central heat laws on $S_N^+$ have rates growing at most linearly, as follows, with $D_k$ being the polynomial in $N$ giving $d_k$, so that $D_0=1$, $D_1=N-1$ and $D_{k+1}=(N-2)D_k-D_{k-1}$:

\begin{theorem}
For $N\geq5$, the rates of a central convolution semigroup of states on $C(S_N^+)$, given by $\varphi_t(u^{(k)}_{ij})=\delta_{ij}e^{-tc_k}$, satisfy $0\leq c_k\leq(N-1)c_1D_k'(N)/D_k(N)$. Thus, if $c_1>0$, the sphere sum $\sum_kd_k^2e^{-Ac_k}$ diverges for $0<A\leq2\eta\sqrt{N(N-4)}/((N-1)c_1)$, where $\eta=arccosh((N-2)/2)$.
\end{theorem}

\begin{proof}
Let $X$ be the character of the magic representation, which is a sum of $N$ projections, so that $0\leq X\leq N$, and $\chi_k=D_k(X)$, by the fusion rules. With $T_j$ being the Chebyshev polynomials, we have $D_k(x)=1+2\sum_{j=1}^kT_j((x-2)/2)$, and since $|T_j'|\leq j^2=T_j'(1)$ on $[-1,1]$, and $T_j'$ is positive and increasing on $[1,\infty)$, we obtain $|D_k'(x)|\leq D_k'(N)$ and $D_k(x)\leq D_k(N)$ for $x\in[0,N]$. By functional calculus, we obtain:
$$0\leq1-\frac{D_k(X)}{D_k(N)}\leq(N-X)\frac{D_k'(N)}{D_k(N)}$$

Now by applying $\varphi_t$, and using $\varphi_t(D_k(X))=d_ke^{-tc_k}$ and $\varphi_t(N-X)=(N-1)(1-e^{-tc_1})$, then dividing by $t$, and letting $t\to0$, we obtain the estimate. Finally, we have $D_k(N)=\sinh((k+1/2)\eta)/\sinh(\eta/2)$, so $\log D_k(N)=k\eta+O(1)$ and $D_k'(N)/D_k(N)=k/\sqrt{N(N-4)}+O(1)$, which gives the divergence. We have checked all this with the program {\tt ribery.py}.
\end{proof}

According to Franz, Kula and Skalski \cite[Theorem 10.10]{fks}, the following rates, which attain the above bound with $c_1=2$, come from a positive central convolution semigroup:
$$c_k=2(N-1)\frac{D_k'(N)}{D_k(N)}$$

At $N=4$ these rates are exactly $k(k+1)$, as checked by {\tt ribery.py}. Thus a positive replacement of the Casimir weights exists, the price to pay being the divergence of the sphere sum at small area, the surface amplitudes and their gluing remaining to be defined. For these rates, logarithmic differentiation of the formula of $D_k(N)$ gives, with $\Delta=\sqrt{N(N-4)}$:
$$\Delta\frac{D_k'(N)}{D_k(N)}=\left(k+\frac{1}{2}\right)\coth\left(\left(k+\frac{1}{2}\right)\eta\right)-\frac{1}{2}\coth\left(\frac{\eta}{2}\right)$$

Thus, at the threshold $A_*=2\eta\Delta/((N-1)c_1)$, the summands of the sphere sum tend to $\exp(\eta\sqrt{N/(N-4)})/(N-4)$, while for $A>A_*$ they decay exponentially, so the threshold in Theorem 10.3 is sharp, as checked with exact integer arithmetic by {\tt ribery.py}. For related work on these processes, see Freslon, Teyssier and Wang \cite{ftw}.

\section{Further questions and outlook}

We have seen in this paper that the complex algebra $B_F$ generated by the Chamseddine--Connes algebra is a twisted version of 15 points, and that the color symmetry of its finite geometry is a twist of a quantum permutation group on 9 points, namely $S_9^+$ in the complex setting, and the quantum group $G_3$ preserving the oriented areas in the real setting. With the trace coming from the fermions, the coaction preserves the weak and color summands separately, and the untwisting turns the matrix algebras into algebras of functions on points, the flip into a graded flip, for $G_3$ in the color case, and the Lie algebras $gl_2,gl_3$ into color Lie algebras.

All this raises a number of questions, and we first list 3 of them, which are quite concrete, and can be investigated with the methods of the present paper:

\begin{enumerate}
\item Can one construct a real spectral triple of KO-dimension $6$, satisfying the first order condition, with algebra $\mathbb C\oplus\mathbb H\oplus\mathbb C^9$, by changing the module data of $F$, for instance via an independent representation, or a module category construction, and how would it be related to $F$? Observe that the monoidal transport of the Heisenberg action and the commuting lifts are excluded, as explained in section 6, and that, by Proposition 6.4, the usual quark module with connected CKM mixing allows at most 3 color points.

\item What are the intermediate quantum groups in Theorem 5.4, in the composite exponent case? Also, since the incidence graph of the Hesse configuration has only classical symmetries, by Theorem 4.8, and adding intertwiners can only shrink this group, which replacements for the ordinary incidence, such as phase-valued tensors on the 9 points, lead to genuinely quantum groups, see \cite{far} for the quantum symmetries of hypergraphs? And since $G_3$ cannot act on the 4 parallel classes of lines, by Proposition 9.5, is there a weaker structure, such as a quantum space of rays, on which the color and weak quantum groups can meet?

\item What happens in the extensions of the Standard Model with a larger algebra, such as the Pati--Salam model \cite{psa}, with algebra $\mathbb H\oplus\mathbb H\oplus M_4(\mathbb C)$ in the spectral version \cite{ccs}, or the algebra $M_2(\mathbb H)\oplus M_4(\mathbb C)$ from \cite{cc1}? Note that at $N=4$ there are 2 choices for $H$, namely $\mathbb Z_4$ and $\mathbb Z_2^2$, leading to different 16-point geometries, distinguished by maximality, by Theorem 5.4, and a selection criterion between them should come from the fermions, the real structure, or the Dirac operator. Also, how does this relate to the twisted spectral triple approach to the breaking of the grand symmetry from \cite{dma}?
\end{enumerate}

Regarding the first question, Proposition 6.4 shows that, as long as the usual quark module, the first order condition and the observed connected CKM mixing are kept, at most 3 color points can act, the bound being attained by the diagonal color algebra. Thus a faithful 9-point replacement must change at least one of the module, support, mass or first order hypotheses of Proposition 6.4, and with the measured masses and mixing fixed, a change in the module data or in the axioms, such as the twisted first order condition from \cite{dma}, is the natural question, the necessary dimension of a replacement module being open. Regarding the second question, Theorems 4.7 and 4.8 show that the quantum symmetries of the 9 points are carried by the orientation of the areas, so the natural candidates are the replacements of the ordinary incidence by structures with phases. As for the third question, the case $N=4$ is the first one where 2 different groups $H$ appear, and is a natural test case for the role of the fermions in the choice of $H$.

Besides the above concrete questions, there are a number of more advanced questions, regarding the passage to the spacetime manifold, the gauge theory, the renormalization and the mixing matrices, as follows:

\begin{enumerate}
\item The passage from the finite geometry to the full Standard Model, on $M\times F$, raises the question of local quantum symmetries, discussed in \cite{bddd}, which is open. By \cite{gsw}, for a smooth action preserving the manifold factor, the action on this factor is classical, although this alone does not classify the actions on the product. Related questions are those of the spectral action prescription from section 8, beyond the classical 1-forms, and of making sense of the untwisted Yang--Mills theory, with the color Lie algebra from Proposition 7.3, and the graded flip from Proposition 7.2.

\item Develop the Yang--Mills program from section 10: the positive replacements of the Casimir weights, from \cite{fks}, and their surface amplitudes, the $S_N^+$ Wilson loops in 2 dimensions and their large $N$ behavior, the lattice theory with magic unitaries on the links, and the relation between the noncrossing combinatorics of $S_{N^2}^+$ and the planar limit of 't Hooft.

\item Regarding renormalization, recall that this is governed by the Hopf algebra of Feynman graphs of Connes and Kreimer \cite{ck1,ck2}, with the Ward--Takahashi identities of QED being implemented as relations in this Hopf algebra, as shown by van Suijlekom \cite{vs2}, and with the underlying universal symmetry being the cosmic Galois group $U^*=U\rtimes\mathbb G_m$ of Connes and Marcolli \cite{cm}, whose normal part $U$ is pro-unipotent, and so of Lie type. A quantum version of $U$, which would provide a link with the present framework, would require a theory of quantum groups beyond the locally compact setting of Kustermans and Vaes \cite{kva}.

\item Regarding the mixing matrices, what is the quantum group invariant measuring the deviation from the Hadamard regime, refining the classical quantity $\delta$ from section 9? Also, what is the role of the quantum subgroups of $S_4^+$, discussed in section 9, in connection with the flavor groups?
\end{enumerate}

Let us also state clearly what the present paper does not do. The algebra $A_F$, its fermion representation, and the mass and mixing data are inputs here, and rewriting them in terms of quantum permutations cannot by itself predict the number of colors or generations, the coupling constants, or the mixing angles. Also, we have no action principle for untwisted quarks, and no new observable or phenomenological fit. What we have is a dictionary, which transports the algebraic structures of the finite geometry to the side of the quantum permutations, and which identifies precisely where this transport works, as for the algebra, its traces and the gauge fields, and where it stops, as for the individual quark triplets, the unoriented areas, or the spectral action with coefficientwise reality.

The main new results here are Theorems 3.4, 4.3, 4.5, 4.7, 4.8, 5.3, 5.4, 8.2, 10.2, 10.3 and Propositions 4.6, 5.2, 6.2, 6.4, 8.3, 9.5, the other statements being either known, as indicated in the text, or of a more heuristic nature.

From a mathematical point of view, the main phenomenon found in this paper is the contrast between the oriented and unoriented area tensors. On the same 9 points, the complex oriented area tensor has an infinite quantum symmetry group, with 216 classical points, while its real part has only classical symmetries, forming a group of order 432, the mechanism being that the adjoint of an intertwining relation reverses the order of the products, while the entrywise conjugation does not. We believe that this phenomenon, where a finite geometry decorated with phases has genuinely quantum symmetries, which disappear when the phases are forgotten, should appear in many other situations. On the physics side, the physical trace and the spectral action provide constraints which are invisible at the level of the fusion rules, the former depending on the fermion representation, and the latter on the reality condition, and we hope that such constraints will help in understanding the untwisting of the whole finite geometry.

These questions are of course not independent, and all depend on a good understanding of the untwisting procedure. As for the order in which all this should be done, QED comes before QCD, and the weak sector, which is classical in disguise, should be understood first, with the Pauli quantum group \cite{bcp} and its quantum subgroups \cite{bbi} as main tools, the color sector, which is genuinely free, coming next. Also, there is no point in formulating a gauge theory for objects which are still twisted, so the Yang--Mills theory and the renormalization should come at the end, while the mixing matrices, which provide a direct contact with the experimental data, should be investigated in parallel. Also, the questions from the first list above, which are finite dimensional, can be investigated by computer, with the methods of the present paper, while those from the second list will certainly require new tools. Thus, a natural order for this program is as follows:

\medskip

\noindent {\bf Step 1.} The electroweak sector, where the algebra is $\mathbb C\oplus\mathbb C\oplus M_2(\mathbb C)$, corresponding to $1+1+4$ points, and where the untwisting is governed by $S_4^+$, the Pauli matrices and a graded flip given by signs. This is the simplest case, and the subgroups of $S_4^+$, which are explicitly known, should allow a detailed analysis, including the role of the circle surviving from the weak symmetry in the computation of Bhowmick, D'Andrea and Dabrowski, and of the Higgs field, as an inner fluctuation.

\medskip

\noindent {\bf Step 2.} The color sector, with the untwisting of $M_3(\mathbb C)$ into $9$ points, the quarks becoming modules over the twisted group algebra of $\mathbb Z_3^2$, and with the color Lie algebra from section 7 replacing $gl_3$. Here the main problem is that of formulating the fermionic part of the theory, compatibly with the real structure and with the first order condition, which requires changing the module data, the monoidal transport of the Heisenberg action and the commuting lifts being both excluded, as explained in section 6.

\medskip

\noindent {\bf Step 3.} The passage from the finite geometry $F$ to the full product $M\times F$, with the inner fluctuations, the spectral action, and the gauge fields written in terms of the untwisted objects, and with the question of local quantum symmetries.

\medskip

\noindent {\bf Step 4.} The Yang--Mills theory, with positive replacements of the Casimir weights, and the renormalization, as discussed in section 10 and in the above questions, which should normally come at the end, once the whole Standard Model is untwisted.

\medskip


Summarizing, the present paper is rather an outline of a program, and we intend to develop it in a series of forthcoming papers.

\begin{thebibliography}{99}

\bibitem{afe} G. Altarelli and F. Feruglio, Tri-bimaximal neutrino mixing from discrete symmetry in extra dimensions, {\em Nuclear Phys. B} {\bf 720} (2005), 64--88.

\bibitem{ado} M. Artebani and I. Dolgachev, The Hesse pencil of plane cubic curves, {\em Enseign. Math.} {\bf 55} (2009), 235--273.

\bibitem{dby} F.P. An et al., Observation of electron-antineutrino disappearance at Daya Bay, {\em Phys. Rev. Lett.} {\bf 108} (2012), 171803.

\bibitem{ase} M. Anshelevich and A.N. Sengupta, Quantum free Yang--Mills on the plane, {\em J. Geom. Phys.} {\bf 62} (2012), 330--343.

\bibitem{ban} T. Banica, Symmetries of a generic coaction, {\em Math. Ann.} {\bf 314} (1999), 763--780.

\bibitem{bhm} T. Banica, Quantum permutations, Hadamard matrices, and the search for matrix models, {\em Banach Center Publ.} {\bf 98} (2012), 11--42.

\bibitem{bbi} T. Banica and J. Bichon, Quantum groups acting on 4 points, {\em J. Reine Angew. Math.} {\bf 626} (2009), 75--114.

\bibitem{bbc} T. Banica, J. Bichon and S. Curran, Quantum automorphisms of twisted group algebras and free hypergeometric laws, {\em Proc. Amer. Math. Soc.} {\bf 139} (2011), 3961--3971.

\bibitem{bbn} T. Banica, J. Bichon and S. Natale, Finite quantum groups and quantum permutation groups, {\em Adv. Math.} {\bf 229} (2012), 3320--3338.

\bibitem{bcp} T. Banica and B. Collins, Integration over the Pauli quantum group, {\em J. Geom. Phys.} {\bf 58} (2008), 942--961.

\bibitem{bne} T. Banica and I. Nechita, Flat matrix models for quantum permutation groups, {\em Adv. in Appl. Math.} {\bf 83} (2017), 24--46.

\bibitem{bdd} J. Bhowmick, F. D'Andrea and L. Dabrowski, Quantum isometries of the finite noncommutative geometry of the Standard Model, {\em Comm. Math. Phys.} {\bf 307} (2011), 101--131.

\bibitem{bddd} J. Bhowmick, F. D'Andrea, B. Das and L. Dabrowski, Quantum gauge symmetries in noncommutative geometry, {\em J. Noncommut. Geom.} {\bf 8} (2014), 433--471.

\bibitem{bgo} J. Bhowmick and D. Goswami, Quantum isometry groups: examples and computations, {\em Comm. Math. Phys.} {\bf 285} (2009), 421--444.

\bibitem{bia} P. Biane, Free Brownian motion, free stochastic calculus and random matrices, {\em Fields Inst. Commun.} {\bf 12} (1997), 1--19.

\bibitem{by} J. Bichon and R. Yuncken, Quantum subgroups of the compact quantum group $SU_{-1}(3)$, {\em Bull. Lond. Math. Soc.} {\bf 46} (2014), 315--328.

\bibitem{bgu} G. Borot and A. Guionnet, Asymptotic expansion of $\beta$ matrix models in the one-cut regime, {\em Comm. Math. Phys.} {\bf 317} (2013), 447--483.

\bibitem{behy} M. Brannan, F. Elzinga, S.J. Harris and M. Yamashita, Crossed product equivalence of quantum automorphism groups of finite dimensional $C^*$-algebras, {\em Int. Math. Res. Not. IMRN} {\bf 2023} (2023), no. 20, 17749--17787.

\bibitem{bgg} G.C. Branco, J.M. G\'erard and W. Grimus, Geometrical T-violation, {\em Phys. Lett. B} {\bf 136} (1984), 383--386.

\bibitem{bcds} T. Brzezi\'nski, N. Ciccoli, L. Dabrowski and A. Sitarz, Twisted reality condition for Dirac operators, {\em Math. Phys. Anal. Geom.} {\bf 19} (2016), Art. 16.

\bibitem{cab} N. Cabibbo, Time reversal violation in neutrino oscillation, {\em Phys. Lett. B} {\bf 72} (1978), 333--335.

\bibitem{cc1} A.H. Chamseddine and A. Connes, Why the Standard Model, {\em J. Geom. Phys.} {\bf 58} (2008), 38--47.

\bibitem{cc2} A.H. Chamseddine and A. Connes, The spectral action principle, {\em Comm. Math. Phys.} {\bf 186} (1997), 731--750.

\bibitem{ccm} A.H. Chamseddine, A. Connes and M. Marcolli, Gravity and the standard model with neutrino mixing, {\em Adv. Theor. Math. Phys.} {\bf 11} (2007), 991--1089.

\bibitem{ccs} A.H. Chamseddine, A. Connes and W.D. van Suijlekom, Beyond the spectral Standard Model: emergence of Pati--Salam unification, {\em J. High Energy Phys.} {\bf 11} (2013), 132.

\bibitem{cpv} M.C. Chen, M. Fallbacher, K.T. Mahanthappa, M. Ratz and A. Trautner, CP violation from finite groups, {\em Nuclear Phys. B} {\bf 883} (2014), 267--305.

\bibitem{cte} Claude and Teo, On the maximality conjecture for quantum permutation groups, Hexagon preprint 2611.00001 (2026).

\bibitem{con} A. Connes, Noncommutative geometry and the standard model with neutrino mixing, {\em J. High Energy Phys.} {\bf 11} (2006), 081.

\bibitem{co2} A. Connes, A unitary invariant in Riemannian geometry, {\em Int. J. Geom. Methods Mod. Phys.} {\bf 5} (2008), 1215--1242.

\bibitem{co3} A. Connes, On the spectral characterization of manifolds, {\em J. Noncommut. Geom.} {\bf 7} (2013), 1--82.

\bibitem{ck1} A. Connes and D. Kreimer, Hopf algebras, renormalization and noncommutative geometry, {\em Comm. Math. Phys.} {\bf 199} (1998), 203--242.

\bibitem{ck2} A. Connes and D. Kreimer, Renormalization in quantum field theory and the Riemann--Hilbert problem I: the Hopf algebra structure of graphs and the main theorem, {\em Comm. Math. Phys.} {\bf 210} (2000), 249--273.

\bibitem{cla} A. Connes and G. Landi, Noncommutative manifolds, the instanton algebra and isospectral deformations, {\em Comm. Math. Phys.} {\bf 221} (2001), 141--159.

\bibitem{cm} A. Connes and M. Marcolli, Noncommutative geometry, quantum fields and motives, AMS (2008).

\bibitem{cmo} A. Connes and H. Moscovici, Type III and spectral triples, in ``Traces in number theory, geometry and quantum fields'', Vieweg (2008), 57--71.

\bibitem{dl1} A. Dahlqvist and T. Lemoine, Large $N$ limit of Yang--Mills partition function and Wilson loops on compact surfaces, {\em Probab. Math. Phys.} {\bf 4} (2023), 849--890.

\bibitem{dl2} A. Dahlqvist and T. Lemoine, Large $N$ limit of the Yang--Mills measure on compact surfaces II: Makeenko--Migdal equations and the planar master field, {\em Forum Math. Sigma} {\bf 13} (2025), e17.

\bibitem{dda} F. D'Andrea and L. Dabrowski, The Standard Model in noncommutative geometry and Morita equivalence, {\em J. Noncommut. Geom.} {\bf 10} (2016), 551--578.

\bibitem{dvv} A. De Rijdt and N. Vander Vennet, Actions of monoidally equivalent compact quantum groups and applications to probabilistic boundaries, {\em Ann. Inst. Fourier} {\bf 60} (2010), 169--216.

\bibitem{dma} A. Devastato and P. Martinetti, Twisted spectral triple for the Standard Model and spontaneous breaking of the grand symmetry, {\em Math. Phys. Anal. Geom.} {\bf 20} (2017), 1--43.

\bibitem{doi} Y. Doi, Braided bialgebras and quadratic algebras, {\em Comm. Algebra} {\bf 21} (1993), 1731--1749.

\bibitem{dou} M.R. Douglas, D-branes and discrete torsion, {\tt arxiv:hep-th/9807235}.

\bibitem{est} L.L. Everett and A.J. Stuart, Icosahedral ($A_5$) family symmetry and the golden ratio prediction for solar neutrino mixing, {\em Phys. Rev. D} {\bf 79} (2009), 085005.

\bibitem{ffz} D.B. Fairlie, P. Fletcher and C.K. Zachos, Trigonometric structure constants for new infinite-dimensional algebras, {\em Phys. Lett. B} {\bf 218} (1989), 203--206.

\bibitem{fma} M. Filaci and P. Martinetti, A critical survey of twisted spectral triples beyond the Standard Model, {\em J. Phys. A} {\bf 56} (2023), 153001.

\bibitem{far} N. Faro\ss, Quantum automorphism groups of hypergraphs, {\em Canad. Math. Commun.} {\bf 1} (2026), e5.

\bibitem{fks} U. Franz, A. Kula and A. Skalski, L\'evy processes on quantum permutation groups, in ``Noncommutative analysis, operator theory and applications'', Springer (2016), 193--259.

\bibitem{ftw} A. Freslon, L. Teyssier and S. Wang, Cutoff profiles for quantum L\'evy processes and quantum random transpositions, {\em Probab. Theory Related Fields} {\bf 183} (2022), 1285--1327.

\bibitem{gal} C. Galindo, Splitting of Clifford groups associated to finite abelian groups, {\tt arxiv:2603.24743}.

\bibitem{gos} D. Goswami, Quantum group of isometries in classical and noncommutative geometry, {\em Comm. Math. Phys.} {\bf 285} (2009), 141--160.

\bibitem{gsw} D. Goswami, Non-existence of genuine (compact) quantum symmetries of compact, connected smooth manifolds, {\em Adv. Math.} {\bf 369} (2020), 107181.

\bibitem{gre} O.W. Greenberg, Spin and unitary-spin independence in a paraquark model of baryons and mesons, {\em Phys. Rev. Lett.} {\bf 13} (1964), 598--602.

\bibitem{gll} W. Grimus, L. Lavoura and P.O. Ludl, Is $S_4$ the horizontal symmetry of tri-bimaximal lepton mixing?, {\em J. Phys. G} {\bf 36} (2009), 115007.

\bibitem{haa} U. Haagerup, Orthogonal maximal abelian $*$-subalgebras of the $n\times n$ matrices and cyclic $n$-roots, in ``Operator algebras and quantum field theory'', International Press (1997), 296--322.

\bibitem{hps} P.F. Harrison, D.H. Perkins and W.G. Scott, Tri-bimaximal mixing and the neutrino oscillation data, {\em Phys. Lett. B} {\bf 530} (2002), 167--173.

\bibitem{hls} M. Holthausen, M. Lindner and M.A. Schmidt, CP and discrete flavour symmetries, {\em J. High Energy Phys.} {\bf 4} (2013), 122.

\bibitem{ivi} I.D. Ivanovi\'c, Geometrical description of quantal state determination, {\em J. Phys. A} {\bf 14} (1981), 3241--3245.

\bibitem{jun} JUNO Collaboration, Measurement of reactor neutrino oscillation with the first JUNO data, {\em Nature} {\bf 654} (2026), 343--348.

\bibitem{ksw} P. Kasprzak, P.M. So\l tan and S.L. Woronowicz, Quantum automorphism groups of finite quantum groups are classical, {\em J. Geom. Phys.} {\bf 89} (2015), 32--37.

\bibitem{kra} T. Krajewski, Classification of finite spectral triples, {\em J. Geom. Phys.} {\bf 28} (1998), 1--30.

\bibitem{kva} J. Kustermans and S. Vaes, Locally compact quantum groups, {\em Ann. Sci. \'Ecole Norm. Sup.} {\bf 33} (2000), 837--934.

\bibitem{lam} C.S. Lam, Determining horizontal symmetry from neutrino mixing, {\em Phys. Rev. Lett.} {\bf 101} (2008), 121602.

\bibitem{lma} G. Landi and P. Martinetti, On twisting real spectral triples by algebra automorphisms, {\em Lett. Math. Phys.} {\bf 106} (2016), 1499--1530.

\bibitem{lem} T. Lemoine, Large $N$ behaviour of the two-dimensional Yang--Mills partition function, {\em Combin. Probab. Comput.} {\bf 31} (2022), 144--165.

\bibitem{lev} T. L\'evy, The master field on the plane, {\em Ast\'erisque} {\bf 388} (2017).

\bibitem{mra} E. Ma and G. Rajasekaran, Softly broken $A_4$ symmetry for nearly degenerate neutrino masses, {\em Phys. Rev. D} {\bf 64} (2001), 113012.

\bibitem{maj} S. Majid, Foundations of quantum group theory, Cambridge Univ. Press (1995).

\bibitem{mvs} M. Marcolli and W.D. van Suijlekom, Gauge networks in noncommutative geometry, {\em J. Geom. Phys.} {\bf 75} (2014), 71--91.

\bibitem{msi} P. Martinetti and D. Singh, Lorentzian fermionic action by twisting Euclidean spectral triples, {\em J. Noncommut. Geom.} {\bf 16} (2022), 513--559.

\bibitem{mig} A.A. Migdal, Recursion equations in gauge field theories, {\em Sov. Phys. JETP} {\bf 42} (1975), 413--418.

\bibitem{nrs} G. Nebe, E.M. Rains and N.J.A. Sloane, The invariants of the Clifford groups, {\em Des. Codes Cryptogr.} {\bf 24} (2001), 99--121.

\bibitem{nr2} G. Nebe, E.M. Rains and N.J.A. Sloane, Self-dual codes and invariant theory, Springer (2006).

\bibitem{psi} M. Paschke and A. Sitarz, Discrete spectral triples and their symmetries, {\em J. Math. Phys.} {\bf 39} (1998), 6191--6205.

\bibitem{pz} J. Patera and H. Zassenhaus, The Pauli matrices in $n$ dimensions and finest gradings of simple Lie algebras of type $A_{n-1}$, {\em J. Math. Phys.} {\bf 29} (1988), 665--673.

\bibitem{psa} J.C. Pati and A. Salam, Lepton number as the fourth ``color'', {\em Phys. Rev. D} {\bf 10} (1974), 275--289.

\bibitem{pop} S. Popa, Orthogonal pairs of $*$-subalgebras in finite von Neumann algebras, {\em J. Operator Theory} {\bf 9} (1983), 253--268.

\bibitem{rwy} V. Rittenberg and D. Wyler, Generalized superalgebras, {\em Nuclear Phys. B} {\bf 139} (1978), 189.

\bibitem{rw2} V. Rittenberg and D. Wyler, Sequences of $\mathbb Z_2\otimes\mathbb Z_2$ graded Lie algebras and superalgebras, {\em J. Math. Phys.} {\bf 19} (1978), 2193.

\bibitem{sch} M. Scheunert, Generalized Lie algebras, {\em J. Math. Phys.} {\bf 20} (1979), 712--720.

\bibitem{sol} P.M. So\l tan, Quantum families of maps and quantum semigroups on finite quantum spaces, {\em J. Geom. Phys.} {\bf 59} (2009), 354--368.

\bibitem{tzy} W. Tadej and K. \.Zyczkowski, A concise guide to complex Hadamard matrices, {\em Open Syst. Inf. Dyn.} {\bf 13} (2006), 133--177.

\bibitem{tho} G. 't Hooft, A planar diagram theory for strong interactions, {\em Nucl. Phys. B} {\bf 72} (1974), 461--473.

\bibitem{vaf} C. Vafa, Modular invariance and discrete torsion on orbifolds, {\em Nuclear Phys. B} {\bf 273} (1986), 592--606.

\bibitem{vwi} C. Vafa and E. Witten, On orbifolds with discrete torsion, {\em J. Geom. Phys.} {\bf 15} (1995), 189--214.

\bibitem{vsu} W.D. van Suijlekom, Noncommutative geometry and particle physics, Springer (2015).

\bibitem{vs2} W.D. van Suijlekom, The Hopf algebra of Feynman graphs in quantum electrodynamics, {\em Lett. Math. Phys.} {\bf 77} (2006), 265--281.

\bibitem{wan} S. Wang, Quantum symmetry groups of finite spaces, {\em Comm. Math. Phys.} {\bf 195} (1998), 195--211.

\bibitem{wil} K.G. Wilson, Confinement of quarks, {\em Phys. Rev. D} {\bf 10} (1974), 2445--2459.

\bibitem{wit} E. Witten, On quantum gauge theories in two dimensions, {\em Comm. Math. Phys.} {\bf 141} (1991), 153--209.

\bibitem{wi2} E. Witten, Two dimensional gauge theories revisited, {\em J. Geom. Phys.} {\bf 9} (1992), 303--368.

\bibitem{wfi} W.K. Wootters and B.D. Fields, Optimal state-determination by mutually unbiased measurements, {\em Ann. Physics} {\bf 191} (1989), 363--381.

\bibitem{wol} L. Wolfenstein, Oscillations among three neutrino types and CP violation, {\em Phys. Rev. D} {\bf 18} (1978), 958--960.

\end{thebibliography}

\end{document}
